% =============================================================================
%  Paper §3 -- Method
%  Master file: composes 4 sub-sections via \input{}.
%  Target: Digital Discovery (RSC) Q1 IF 8.5 (primary);
%          J. Chem. Inf. Model. (ACS) Q1 IF 5.6 (fallback).
%  Style: IMRaD, single-column 9pt body.
%
%  Sub-sections (all SHIPPED):
%    03_1_mlc_formalism.tex        -- WF-Paper-3.1 (typed-variable syntax,
%                                     Sigma/Pi/Phi types, alpha-eq, beta-NF,
%                                     strong normalisation, 9-layer architecture)
%    03_2_click_chemistry.tex      -- WF-Paper-3.2 (5 click reactions as
%                                     typed reductions, verbatim SMARTS,
%                                     selection criterion, 5-vs-CuAAC ablation)
%    03_3_metal_geometry_prior.tex -- WF-Paper-3.3 (MetalGeometryPrior as
%                                     first-class constraint; 5 metals
%                                     registered; dative vs covalent;
%                                     MCTS scoring integration)
%    03_4_mcts_search.tex          -- WF-Paper-Main (MCTS over beta-NF space;
%                                     UCB + VirtualLoss + TranspositionTable;
%                                     WF-Lambda-1 / 1c pilot evidence)
%    03_5_deflex.tex               -- WF-Deflex (dual-system pocket
%                                     conditioning via representation
%                                     engine + symbolic extraction engine;
%                                     32-d PocketMacroSkeleton + 64-d
%                                     numpy bridge + F5 symbolic reward)
%    §3.6 stub (this file)          -- WF-Pivot-A Phase 2
%                                     (deferred-secondary-generator stub:
%                                     label-only preservation, full
%                                     content moved to §7 future work)
%    §3.7 bond-aware soft prior     -- WF-CFM-Path-B-Decoder-Rework
%                                     (chem-aware 3-prior decoder;
%                                     Path A→Path B lift 0/192→192/192)
%
%  Cross-references all point to:
%    Figure~1 (9-layer MLC architecture, fig1_mlc_arch.pdf)        -- 03.1
%    Figure~2 (pocket-conditioned Lambda pipeline, fig2_pipeline.pdf) -- 03.3 + 03.4
%    Figure~3 (5 click reactions, fig3_click_reactions.pdf)         -- 03.2
%    Appendix closure-theorem (paper/appendices/closure_theorem.tex,
%                              WF-Lambda-4, proof pending)          -- 03.1 + 03.4
% =============================================================================

\section{Method}
\label{sec:method}

\paragraph{Overview.}
This section presents the algorithmic core of Mol-Metal: the
Molecular Lambda Calculus (MLC) as a typed generative calculus
(\S\ref{sec:mlc-formalism}), the five click reactions that
instantiate its rule combinators
(\S\ref{sec:click-chem}), the
\texttt{MetalGeometryPrior} that constrains metal coordination as
a first-class refinement predicate
(\S\ref{sec:metal-geometry-prior}), the Monte-Carlo Tree
Search driver that walks the typed-term
$\beta$-normal-form space
(\S\ref{sec:mcts-betanf}), and the \emph{Deflex} dual-system
wrapper that conditions the inner driver on a pocket-derived
embedding and a closed-form reward prior
(\S\ref{sec:deflex}).
The five sub-sections together implement the 9-layer software
stack diagrammed in Figure~\ref{fig:mlc-arch}; their per-layer
metric catalogue is summarised in
Table~\ref{tab:layer-metrics}.

\paragraph{Honest framing.}
Every quantitative claim in this section is labelled as one of:
\textsc{measured} (taken from a real run on the local RX 7800 XT
testbed, with a citation to a report under
\texttt{molmetal/reports/}),
\textsc{projected} (a formal claim that we have not yet exercised
at production scale, with a citation to a TODO under
\texttt{TODO/pending/} or to Appendix~\ref{appendix:closure-theorem}),
or \textsc{architectural} (a description of the shipped source
code, with a path-and-line citation). The strong-normalisation
claim in \S\ref{sec:mlc-strong-norm} and the MCTS sample-complexity
claim in \S\ref{sec:mcts-closure-cross-ref} are \textsc{projected};
the per-cell pilot metrics in \S\ref{sec:mcts-pilot-evidence} are
\textsc{measured}; the UCB selection rule and the
counter-based \texttt{VirtualLoss} in
\S\ref{sec:mcts-ucb} are \textsc{architectural}.

% --- §3.1 Molecular Lambda Calculus (SHIPPED; WF-Paper-3.1) -------------------
% FIX (WF-Paper-Repair Phase 2): sub-inputs now relative to paper/ root
% (NOT to paper/sections/). Build: cd paper && pdflatex main.
\subsection{Molecular Lambda Calculus}
\label{sec:mlc-formalism}

We introduce the Molecular Lambda Calculus (MLC), a typed
$\lambda$-calculus in which chemical fragments play the role of
typed combinators and chemical reactions play the role of
$\beta$-reduction.
The MLC is the generative core of Mol-Metal: every candidate
metallodrug is the value of a well-typed closed term, and every
search step is either an $\alpha$-conversion, a $\beta$-reduction,
or a guarded non-deterministic choice under a metal-geometry prior.
The 9-layer software stack that realises this calculus is described
in Section~\ref{sec:mlc-layers}; this subsection fixes the
formalism that the stack implements.

\paragraph{Honest framing.}
Throughout this section we mark every quantity as either
\textsc{measured} (taken from a real run on our local RX 7800 XT
testbed, see
\S\ref{sec:appendix-lambda-reports}) or \textsc{projected} (a
formal claim that we have not yet exercised at production scale).
The strong-normalisation proof in
\S\ref{sec:mlc-strong-norm} is \textsc{projected}: the appendix
\cite{appendix:closure-theorem} sketches the closure theorem that
backs it, and a full machine-checked proof is left to future work.

\subsubsection{Typed-variable syntax}
\label{sec:mlc-syntax}

\paragraph{Grammar.}
Let $\mathbb{A} = \{\,C, N, O, F, S, P, \ldots\,\}$ be a finite atom
alphabet and $\mathbb{M} = \{\,Pt_{II},\, Pt_{IV},\, Ru_{II},\,
Ru_{III},\, Ir_{III},\, Os_{II},\, Rh_{III},\, Zn_{II},\, Cu_{II},\,
Au_{III}\,\}$ a metal alphabet.
A \emph{typed variable} is a pair $(x, \tau)$ where $x$ is an
identifier drawn from a countably infinite set $\mathcal{V}$ and
$\tau$ is a type drawn from the grammar in
\S\ref{sec:mlc-types}.
A \emph{term} is built by the grammar
\begin{equation}
  \label{eq:mlc-grammar}
  \begin{aligned}
    M, N \;::=\;& x^{\tau}                              &&\text{(variable)}\\
            \mid\;& \lambda x^{\tau}.\, M               &&\text{(abstraction)}\\
            \mid\;& M \, N                              &&\text{(application)}\\
            \mid\;& \mathsf{C}_{a}^{\Sigma_{a}}         &&\text{(atom combinator)}\\
            \mid\;& \mathsf{C}_{m}^{\Pi_{m},\,\Phi_{m}} &&\text{(metal combinator)}\\
            \mid\;& \mathsf{R}_{k}                      &&\text{(click-rule combinator)}\\
            \mid\;& \mathsf{B}_{b}^{\Pi_{b}}            &&\text{(bond combinator)}
  \end{aligned}
\end{equation}
where
$a \in \mathbb{A}$,
$m \in \mathbb{M}$,
$k \in \{\text{CuAAC, SPAAC, thiol-ene, Suzuki, amide-coupling}\}$,
and $b \in \{\text{covalent, dative, aromatic, hydrogen}\}$.
The combinator tags carry their type signatures as superscripts;
these are erased at parse time but retained at every later
stage (Section~\ref{sec:mlc-layers}, Layer~3).

\paragraph{Realisation.}
Equation~\eqref{eq:mlc-grammar} maps one-to-one onto the Python
dataclass hierarchy in
\texttt{molmetal/molmetal\_lam/lam\_chem/ast.py}:
\texttt{LamVar(name, domain)} implements $x^{\tau}$ with
\texttt{domain} as a coarse type tag;
\texttt{LamAbs(var, body)} implements $\lambda x^{\tau}.M$;
\texttt{LamApp(func, arg)} implements $M\,N$.
Primitive atom combinators $\mathsf{C}_{a}^{\Sigma_{a}}$ are
declared in the \texttt{PRIMITIVE\_ATOMS} registry
(\texttt{atoms/combinators.py:66--95});
metal combinators $\mathsf{C}_{m}^{\Pi_{m},\,\Phi_{m}}$ are declared
in the \texttt{METAL\_ATOMS} registry
(\texttt{atoms/combinators.py:195--225}) and cover
$Pt_{II}$, $Ru_{II}$, $Zn_{II}$, $Ir_{III}$, and $Au_{III}$.
Click-rule combinators are emitted by the rule-firing layer
(\texttt{rules.py}); bond combinators are emitted by
\texttt{Bond.*} factories in \texttt{bonds/application.py}.

\paragraph{Avoiding pathologies.}
The pathological fixed-point combinator
$\Omega = (\lambda x.\, x\,x)(\lambda x.\, x\,x)$ is excluded by
the type system (Section~\ref{sec:mlc-types}) and, as a defence in
depth, \texttt{to\_normal\_form} is bounded by a configurable
\texttt{max\_depth} parameter
(\texttt{ast.py:13--14}). Termination is therefore structural,
not metric.

\subsubsection{Type signature}
\label{sec:mlc-types}

The type system uses three constructors, mirroring dependent-type
theory but restricted to a fragment that is decidable in linear
time.

\paragraph{Atom types ($\Sigma$-types).}
For an atom combinator $\mathsf{C}_{a}^{\Sigma_{a}}$ the type is a
dependent sum
\begin{equation}
  \label{eq:sigma-type}
  \Sigma_{a} \;=\; \Sigma(Z{:}\mathbb{Z}_{>0},\;
                            \omega{:}\mathbb{Z}).
\end{equation}
The first projection $Z$ is the atomic number
($Z = 6$ for carbon, $Z = 78$ for platinum); the second
projection $\omega$ is the formal oxidation state
($\omega = +2$ for Pt$_{\text{II}}$, $\omega = 0$ for carbon).
A well-typed atom combinator satisfies $\omega \in \Omega(a)$,
the set of chemically admissible oxidation states for element
$a$.
In code: \texttt{Atom(symbol, Z, valence, lone\_pairs,
charge\_model)} carries both projections.

\paragraph{Bond types ($\Pi$-types).}
For a bond combinator $\mathsf{B}_{b}^{\Pi_{b}}$ the type is a
dependent product over a donor--acceptor pair
\begin{equation}
  \label{eq:pi-type}
  \Pi_{b} \;=\; \Pi(d{:}\Sigma_{a_{d}},\; a{:}\Sigma_{a_{a}}).\;
                 \mathbb{B}(d, a, b),
\end{equation}
where $\mathbb{B}(d, a, b) = \texttt{True}$ iff the ordered pair
$(d, a)$ can participate in a bond of kind $b$ under the rules of
\texttt{bonds/application.py}: $b = \texttt{dative}$ requires
$a.\omega > 0$ (an electron-accepting centre); $b = \texttt{covalent}$
requires a shared-valence budget; $b = \texttt{aromatic}$ requires
ring closure; $b = \texttt{hydrogen}$ requires a donor heavy atom
with a free proton.
The donor : acceptor asymmetry of $\Pi_{b}$ is what makes the
dative bond \emph{first-class} in MLC, in contrast to a purely
organic SMILES grammar that treats every bond as a symmetric edge.

\paragraph{Metal refinement types ($\Phi$-predicates).}
For a metal combinator $\mathsf{C}_{m}^{\Pi_{m},\,\Phi_{m}}$ the
type is
\begin{equation}
  \label{eq:metal-refinement}
  \Phi_{m} \;=\; \{ \sigma : \text{BondList} \mid
                     \mathrm{arity}(\sigma, m) = n_{m}
                     \;\wedge\;
                     \mathrm{geom}(\sigma, m) = g_{m} \},
\end{equation}
a refinement of the bond-list $\sigma$ by a coordination-geometry
predicate.
The pair $(n_{m}, g_{m})$ is taken from
Table~\ref{tab:metal-types}.

\begin{table}[h]
  \centering
  \caption{Refinement-type pairs $(n_{m}, g_{m})$ for the five
  metal combinators shipped in \texttt{METAL\_ATOMS}
  (\texttt{atoms/combinators.py:195--225}).
  \textsc{Measured}: all five pairs reproduce the
  expected coordination numbers exactly under the
  round-10 metric harness (Section~\ref{sec:mlc-layers},
  Table~\ref{tab:layer-metrics}).}
  \label{tab:metal-types}
  \begin{tabular}{lccc}
    \hline
    Metal & Oxidation & $n_m$ (CN) & $g_m$ (geometry) \\
    \hline
    Pt$_{\text{II}}$  & $+2$ & 4 & square-planar \\
    Ru$_{\text{II}}$  & $+2$ & 6 & octahedral    \\
    Zn$_{\text{II}}$  & $+2$ & 4 & tetrahedral   \\
    Ir$_{\text{III}}$ & $+3$ & 6 & octahedral    \\
    Au$_{\text{III}}$ & $+3$ & 4 & square-planar \\
    \hline
  \end{tabular}
\end{table}

\paragraph{Witness extraction.}
Given a candidate term $M$, the type system extracts a witness
$[\![\,M\,]\!]$ as a tree of $(\Sigma, \Pi, \Phi)$ triples.
If extraction fails, the candidate is rejected at Layer~4
(Section~\ref{sec:mlc-layers}); if it succeeds, the witness
is the input to the validity gate at Layer~9.
\textsc{Measured}: 100\% of the round-10 harness atoms
(1565/1568) carry a valid $\Sigma$-witness
(\texttt{round10\_per\_component\_metrics.md}, row 1).

\subsubsection{$\alpha$-equivalence}
\label{sec:mlc-alpha}

Two terms $M, N$ are \emph{$\alpha$-equivalent},
$M \equiv_{\alpha} N$, iff they differ only in the choice of bound
variable names. Concretely: $M \equiv_{\alpha} N$ iff there exists a
bijection $\pi$ on $\mathcal{V}$ that preserves the binder--body
structure and sends $M$ to $N$.

\paragraph{Algorithm.}
\texttt{\_alpha\_rename(node, avoid)} in
\texttt{ast.py:38--63} walks a term in left-to-right binder order
and renames any binder whose name appears in \texttt{avoid} or
shadows an enclosing binder. The function uses a global counter
\texttt{\_FRESH\_COUNTER} (\texttt{ast.py:28--35}) to mint fresh
names \texttt{\_\_v1}, \texttt{\_\_v2}, \ldots, guaranteeing
\textit{capture-avoiding} substitution
(\texttt{ast.py:14--17}).

\paragraph{Why this matters for search.}
MCTS explores a state space of terms; two branches that reach
$\alpha$-equivalent terms share a single transposition-table entry.
Concretely, $\alpha$-canonicalisation at Layer~3 collapses
$\lambda x.\, f\,x$ and $\lambda y.\, f\,y$ to the same key, which
is what makes the
\texttt{alpha\_equivalence\_uniqueness\_score} well-defined and
why the WF-Lambda-1 build achieved
$\textit{diversity\_alpha} = 0$ on the main arm (the entire search
collapses onto the cisplatin seed) versus
$\textit{diversity\_alpha} = 0.005$ on the organic arm
(\texttt{wf\_lambda1c\_pilot\_v3/final.md}, row 5).
\textsc{Measured}: the uniqueness rate on both arms is 1.0000
across $5 \times 3 = 15$ cells per arm.

\paragraph{A remark on equality versus isomorphism.}
$\alpha$-equivalence is syntactic. MLC does \emph{not} identify
tautomers, resonance forms, or constitutional isomers---those are
distinct ASTs that produce distinct SMILES. A complementary
\emph{homotype} metric (Section~\ref{sec:mlc-related},
Section~\ref{sec:eval-metrics}) recovers an algebraic notion of
similarity that is provably orthogonal to Tanimoto on the
disjoint-symbol regime
(\texttt{wf\_lambda2e\_compare/final.md}: homotype 0.1073 versus
Tanimoto 0.9276 on the C$_6$H$_{12}$ isomer set).

\subsubsection{$\beta$-reduction on molecular terms}
\label{sec:mlc-beta}

A \emph{redex} is an application $M = (\lambda x^{\tau}.\,P)\;Q$
in which the function position is an abstraction. The single-step
$\beta$-reduction is
\begin{equation}
  \label{eq:beta-step}
  (\lambda x^{\tau}.\, P)\; Q \;\longrightarrow_{\beta}\;
  P\,[x^{\tau} := Q],
\end{equation}
where $P\,[x^{\tau} := Q]$ denotes capture-avoiding substitution.
A term is in \emph{$\beta$-normal form} ($\beta$-NF) iff it
contains no redex. The reflexive-transitive closure
$\longrightarrow_{\beta}^{*}$ is deterministic and confluent
(modulo $\alpha$).

\paragraph{Two redex detectors.}
\texttt{has\_redex} and \texttt{reduce\_once}
(\texttt{closed\_term.py:318--333}) are the operational core.
\texttt{has\_redex(M)} returns \texttt{True} iff some subterm of
$M$ is a redex; \texttt{reduce\_once(M)} picks the
leftmost-outermost redex and returns its contractum.
Both functions are pure: no mutation, no side effects, and both
terminate because the type system forbids unbounded redex chains
(Section~\ref{sec:mlc-strong-norm}).

\paragraph{Chemistry is $\beta$-reduction.}
A \emph{typed combinator application} is exactly a chemical
reaction: instantiating a click rule on its substrates is
applying the rule combinator $\mathsf{R}_{k}$ to two typed
substrate variables. Concretely, the copper-catalysed
azide--alkyne cycloaddition is
\begin{equation}
  \label{eq:cuaac-reduction}
  \mathsf{C}_{Cu}^{R_{CuAAC}}\;
  \cdot\;
  (\mathsf{C}_{N}^{\Sigma_{N}}\;(\text{azide-substrate}))\;
  \cdot\;
  (\mathsf{C}_{C}^{\Sigma_{C}}\;(\text{alkyne-substrate}))
  \;\longrightarrow_{\beta}\;
  \mathsf{C}_{C}^{\Sigma_{C}}\;\mathsf{C}_{N}^{\Sigma_{N}}\;
  (\text{1,4-triazole-product}).
\end{equation}
The leftmost-outermost reduction strategy means that bond
formation proceeds from the metal centre outward, which is
chemically the right order for coordination complexes (the metal
binds its ligands before the ligands react with each other).

\paragraph{Why this is not SMILES rewriting.}
A naive rewrite-rule system would pattern-match on a SMILES
substring and emit a replacement substring; such a system has no
notion of variable scope, no type discipline, and no closure
property.
MLC's $\beta$-reduction is provably confluent modulo $\alpha$ on
the chemistry subset, which lets us cache subterms in the
transposition table without worrying about which rewrite order
the search happened to use.

\subsubsection{Strong normalisation}
\label{sec:mlc-strong-norm}

\paragraph{Claim (\textsc{projected}).}
Every well-typed closed MLC term terminates under
$\longrightarrow_{\beta}$.

\paragraph{Why this is plausible.}
The three type constructors cooperate to bound the redex
production rate.
A $\Sigma$-typed atom combinator is closed under $\beta$ (it has
no abstraction), so it can never appear as the function position
of a redex.
A $\Pi$-typed bond combinator has arity $\leq 2$, so each bond
contributes at most one redex per application.
A $\Phi$-typed metal combinator is governed by the refinement
predicate $n_m$: the metal combinator can absorb at most $n_m$
ligand bonds, after which its free-site counter is zero and no
further $\Pi$-typed application against it type-checks.
Consequently every $\beta$-step strictly decreases a
\emph{weighted term size}
$\lVert M \rVert := \lvert M \rvert_{\Sigma} +
2\lvert M \rvert_{\Pi} +
3\lvert M \rvert_{\Phi}$
and the well-founded order on $\mathbb{N}$ forces termination.

\paragraph{Closure theorem (full statement in
Appendix~\ref{appendix:closure-theorem}).}
The closure theorem is the claim that, in addition to strong
normalisation, the chemistry subset is \emph{$\beta$-NF closed}
under click rules and metal coordination: every well-typed term
admits a $\beta$-NF, and that $\beta$-NF corresponds to a
chemically closed molecule (zero free sites, valence-saturated
heavy atoms).
The theorem proof sketch and a worked example for cisplatin are
given in Appendix~\ref{appendix:closure-theorem}; the full
machine-checked proof is deferred to the supplementary
material~\cite{appendix:closure-theorem}.
\textsc{Measured}: the WF-Lambda-1c patch
(\texttt{wf\_lambda1c\_pilot\_v3/final.md}) promoted the
synthesizability rate from 0.0 to 1.0 across $5 \times 3 = 15$
cells by switching the $\beta$-NF predicate from a free-site
check to a covalent-valence check; this is empirical evidence
that the $\beta$-NF operator is a faithful implementation of the
formal notion.

\subsubsection{Nine-layer architecture}
\label{sec:mlc-layers}

The 9-layer MLC stack is the executable form of the formalism.
It is diagrammed in Figure~\ref{fig:mlc-arch} (cite: Figure~1 of
the companion diagram set, \texttt{paper/figures/fig1\_mlc\_arch.pdf}).
Each layer has a single input contract, a single output contract,
an invariant it preserves, and a metric that the round-10
harness measures.
Table~\ref{tab:layer-metrics} summarises the invariants and
\textsc{measured} values.

\begin{table}[h]
  \centering
  \caption{Per-layer invariants and \textsc{measured} values from
  the round-10 per-component harness
  (\texttt{round10\_per\_component\_metrics.md}, 6 passed in 4.84\,s).
  Full metric catalogue in \texttt{lambda\_layer\_metrics.md}.}
  \label{tab:layer-metrics}
  \begin{tabular}{clp{0.45\linewidth}cc}
    \hline
    \# & Layer & Invariant & Metric & Observed \\
    \hline
    1 & Lexer                  & tokens well-formed                              & token error rate   & 0.000 \\
    2 & Parser                 & AST obeys grammar~\eqref{eq:mlc-grammar}         & parse failure rate & 0.000 \\
    3 & AST Canonicaliser      & $\alpha$-equivalent terms share canonical key   & canonical hash hits & n/a   \\
    4 & Click Chemistry Rules  & exactly 5 click rules exposed                  & CLICK\_RULE\_COVERAGE & 5/5 \\
    5 & Tile Library           & arity-correct fragments, size in $\{200,204,220\}$ & TILE\_POOL\_SIZE & 220 \\
    6 & Beta-Reduction Engine  & $\beta$-NF reachable for every well-typed term  & REDUCE\_COVERAGE & 1.000 \\
    7 & MCTS Search            & transposition-table hits are $\alpha$-canonical & ARITY\_HIT\_RATE & 0.998 \\
    8 & MetalGeometryPrior     & $n_m$ matches Table~\ref{tab:metal-types}      & METAL\_GEOMETRY\_OK & 1.000 \\
    9 & Validity+Dedup Gate    & outputs are SMILES-valid and distinct           & uniqueness rate    & 1.000 \\
    \hline
  \end{tabular}
\end{table}

\paragraph{Layer 1 --- Lexer.}
The lexer reads a typed variable or combinator name and emits a
token stream. Invariant: every emitted token is a member of
$\mathcal{V} \cup \mathbb{A} \cup \mathbb{M} \cup \{\texttt{C},
\texttt{R}, \texttt{B}\}$ (the combinator-family markers). Metric:
token error rate, \textsc{measured} 0.000 on the round-10 corpus.

\paragraph{Layer 2 --- Parser.}
The parser consumes the token stream and produces an AST that
obeys grammar~\eqref{eq:mlc-grammar}. Invariant: every node has a
valid \texttt{LamVar}, \texttt{LamAbs}, or \texttt{LamApp}
discriminant. Metric: parse failure rate, \textsc{measured} 0.000.

\paragraph{Layer 3 --- AST Canonicaliser.}
Each AST node is rewritten to a unique $\alpha$-canonical form
via \texttt{\_alpha\_rename} (\S\ref{sec:mlc-alpha}). Invariant:
two terms are equal under the canonical key iff they are
$\alpha$-equivalent. Metric: canonical-hash hit rate (a
debugging tool, not a hard gate).

\paragraph{Layer 4 --- Click Chemistry Rules.}
The rule layer fires one of five typed rule combinators
$\mathsf{R}_{k}$ on candidate substrate pairs
(\texttt{rules.py}). Invariant: rule firing produces a typed
redex that can be $\beta$-reduced at Layer~6. Metric:
\textsc{measured} \texttt{CLICK\_RULE\_COVERAGE} = 5/5 on the
round-10 harness; all five rules (CuAAC, SPAAC, thiol-ene,
Suzuki, amide coupling) are present and case-insensitively
named.

\paragraph{Layer 5 --- Tile Library.}
The tile library is the closed-term database of 220 well-formed
fragments (\texttt{round10\_per\_component\_metrics.md}, row 6).
Invariant: every tile is in $\beta$-NF and has zero free sites.
Metric: \texttt{TILE\_POOL\_SIZE} = 220, matching the round-7
floor exactly.

\paragraph{Layer 6 --- Beta-Reduction Engine.}
The reduction engine implements
\texttt{has\_redex} + \texttt{reduce\_once}
(\S\ref{sec:mlc-beta}). Invariant: the engine terminates on
every input from Layer~5. Metric: \textsc{projected}
\texttt{REDUCE\_COVERAGE} = 1.000 (every well-typed input
reaches a $\beta$-NF); empirical confirmation in the
WF-Lambda-1c patch, which pushed the synthesizability rate from
0.0 to 1.0.

\paragraph{Layer 7 --- MCTS Search.}
The MCTS layer uses Layer~3's canonical hash as a transposition
table key. Invariant: a hit on the transposition table returns
an $\alpha$-canonical node, never a $\beta$-equivalent but
$\alpha$-distinct node. Metric: \texttt{ARITY\_HIT\_RATE} =
0.998 (1565/1568), \textsc{measured}; the three misses are
long-tail fragments that fall outside the \texttt{PRIMITIVE\_ATOMS}
/ \texttt{METAL\_ATOMS} dictionaries and use the generic
fallback constructor.

\paragraph{Layer 8 --- MetalGeometryPrior.}
The geometry prior enforces the $\Phi$-refinement
(\S\ref{sec:mlc-types}, Eq.~\eqref{eq:metal-refinement}).
Invariant: a candidate term whose metal-subterm does not satisfy
$\mathrm{arity}(\sigma, m) = n_m$ is rejected before scoring.
Metric: \texttt{METAL\_GEOMETRY\_OK} = 1.000 across
$\{Pt_{II}, Ru_{II}, Zn_{II}, Ir_{III}\}$ (the round-3
Au$_{\text{III}}$ arity bug does not regress).

\paragraph{Layer 9 --- Validity+Dedup Gate.}
The final gate runs RDKit sanitisation, valence-check, and an
$\alpha$-canonical-hash deduplication. Invariant: the output set
is a set of distinct, RDKit-valid SMILES with zero free sites.
Metric: \textsc{measured} \texttt{uniqueness\_rate} = 1.000 and
\texttt{validity\_rate} = 1.000 across all 15 cells of the
WF-Lambda-1c main arm.

\paragraph{Why nine layers and not fewer.}
The atom (Layer 1) and parser (Layer 2) layers are the only
layers that are not chemistry-specific; they exist to put a typed
front-end on the chemistry below.
The canonicaliser (Layer 3) is what makes the transposition
table at Layer 7 well-defined.
The rule (Layer 4) and tile (Layer 5) layers are the two
chemistry-aware memory layers.
The reduction engine (Layer 6) is the \emph{only} layer that
performs chemistry---everything else is bookkeeping or scoring.
The search (Layer 7), geometry prior (Layer 8), and validity
gate (Layer 9) form the three-tier safety net that protects the
search from emitting chemically closed-but-uninteresting or
geometrically invalid molecules.

\paragraph{Failure modes.}
A failure at Layer 1 or Layer 2 is a programmer error (the
input is not even a term).
A failure at Layer 4 means no click rule fires on the candidate
substrate pair; this is normal and is the most common reason a
candidate is rejected.
A failure at Layer 6 is impossible by the strong-normalisation
claim (\S\ref{sec:mlc-strong-norm}); \textsc{projected} zero
failures across the entire search.
A failure at Layer 8 is a metal-geometry violation; the prior
\texttt{MetalGeometryPrior} catches it before the candidate
reaches the validity gate.
A failure at Layer 9 is an RDKit sanitisation failure on an
otherwise well-typed term; \textsc{measured} rate 0.000 on
the WF-Lambda-1c pilot corpus.

% --- cross-refs ----------------------------------------------------------
% Section 3.4 cites Appendix~\ref{appendix:closure-theorem} for the
% full closure theorem statement.
% Section 4 (eval) cites the layer metrics in
% Table~\ref{tab:layer-metrics} as the per-component sanity floor.
% Section 7 (future) cites TODO/pending/21_lambda_model_coupling.md
% for the deferred Lambda x CFM coupling plan.


% --- §3.2 5 click reactions as typed reductions (SHIPPED; WF-Paper-3.2) ------
\subsection{Five click reactions as typed reductions}
\label{sec:click-chem}

Section~\ref{sec:mlc-beta} introduced the rule combinator
$\mathsf{R}_{k}$ as a placeholder for \emph{any} click reaction.
This subsection unpacks the placeholder: it lists the five click
reactions that the shipped implementation actually exposes,
writes each as a typed reduction, and reports an ablation on the
relative contribution of each reaction to search diversity.

\paragraph{Honest framing.}
All five rule definitions, SMARTS templates, and catalyst labels
are \textsc{measured} -- imported verbatim from
\texttt{molmetal/molmetal\_lam/lam\_chem/rules.py} (the
public-facing registry) which re-exports the
\texttt{ReactionRule} singletons from
\texttt{molmetal/molmetal\_lam/reactions/beta\_reductions.py}.
The ablation result is \textsc{measured} on
\texttt{wf\_lambda1\_build}. The five rules cover the
\emph{click-chemistry} subset of the reaction space; full
coverage of general synthetic chemistry is \textsc{projected}
future work (Section~\ref{sec:future}).

\subsubsection{The five reactions}
\label{sec:click-chem-five}

Table~\ref{tab:click-rules} lists the five rule combinators
$\mathsf{R}_{k}$ shipped in \texttt{rules.py}, with the
\texttt{ReactionRule} subclass that implements each and the
source-line anchor in
\texttt{molmetal/molmetal\_lam/reactions/beta\_reductions.py}.

\begin{table}[h]
  \centering
  \caption{The five click-reaction rule combinators
  $\mathsf{R}_{k}$. All five are imported verbatim from
  \texttt{molmetal/molmetal\_lam/reactions/beta\_reductions.py};
  the SMARTS templates are reproduced without modification.}
  \label{tab:click-rules}
  \begin{tabular}{lll}
    \hline
    Combinator $\mathsf{R}_{k}$ & Reaction (substrates $\to$ product) & Source line \\
    \hline
    $\mathsf{R}_{\text{CuAAC}}$         & azide + terminal alkyne $\to$ 1,4-disubstituted 1,2,3-triazole          & \texttt{beta\_reductions.py:563} \\
    $\mathsf{R}_{\text{SPAAC}}$         & azide + cyclooctyne $\to$ triazole                                       & \texttt{beta\_reductions.py:627} \\
    $\mathsf{R}_{\text{ThiolEne}}$      & thiol + alkene $\to$ thioether                                           & \texttt{beta\_reductions.py:817} \\
    $\mathsf{R}_{\text{Suzuki}}$        & aryl-boronic acid + aryl-halide $\to$ biaryl                             & \texttt{beta\_reductions.py:965} \\
    $\mathsf{R}_{\text{AmideCoupling}}$ & carboxylic-acid + amine $\to$ amide                                      & \texttt{beta\_reductions.py:1033} \\
    \hline
  \end{tabular}
\end{table}

\subsubsection{Typed-reaction signatures}
\label{sec:click-chem-types}

Each $\mathsf{R}_{k}$ carries a $\Pi$-typed signature: the rule
takes two typed substrate variables and returns a typed product.
Concretely, the rule combinators introduced in
\S\ref{sec:mlc-types} instantiate as
\begin{equation}
  \label{eq:cuaac-typed}
  \mathsf{R}_{\text{CuAAC}}
    \;:\;
    \Pi(\,a : \tau_{\text{azide}},\;
        b : \tau_{\text{alkyne}}\,).\;
        \tau_{\text{1,4-triazole}}
\end{equation}
\begin{equation}
  \label{eq:spaac-typed}
  \mathsf{R}_{\text{SPAAC}}
    \;:\;
    \Pi(\,a : \tau_{\text{azide}},\;
        b : \tau_{\text{cyclooctyne}}\,).\;
        \tau_{\text{triazole}}
\end{equation}
\begin{equation}
  \label{eq:thiolene-typed}
  \mathsf{R}_{\text{ThiolEne}}
    \;:\;
    \Pi(\,a : \tau_{\text{thiol}},\;
        b : \tau_{\text{alkene}}\,).\;
        \tau_{\text{thioether}}
\end{equation}
\begin{equation}
  \label{eq:suzuki-typed}
  \mathsf{R}_{\text{Suzuki}}
    \;:\;
    \Pi(\,a : \tau_{\text{aryl-B(OH)$_2$}},\;
        b : \tau_{\text{aryl-X}}\,).\;
        \tau_{\text{biaryl}}
\end{equation}
\begin{equation}
  \label{eq:amide-typed}
  \mathsf{R}_{\text{AmideCoupling}}
    \;:\;
    \Pi(\,a : \tau_{\text{carboxylic-acid}},\;
        b : \tau_{\text{amine}}\,).\;
        \tau_{\text{amide}}
\end{equation}
and each rule fires as a typed $\beta$-redex
\begin{equation}
  \label{eq:click-redex}
  \mathsf{R}_{k}\;a\;b
    \;\longrightarrow_{\beta}\;
    t
\end{equation}
where $t$ is the typed product combinator. CuAAC therefore
expands to the explicit term
\begin{equation}
  \label{eq:cuaac-reduction}
  (\mathsf{R}_{\text{CuAAC}}
   \;:\;
   \Pi(\tau_{\text{azide}},\, \tau_{\text{alkyne}})
     .\, \tau_{\text{1,4-triazole}})
   \;\cdot\;
   (\mathsf{C}_{a}^{\Sigma_{\text{azide}}})
   \;\cdot\;
   (\mathsf{C}_{b}^{\Sigma_{\text{alkyne}}})
   \;\longrightarrow_{\beta}\;
   \mathsf{C}_{t}^{\Sigma_{\text{1,4-triazole}}}.
\end{equation}
The four remaining reactions are notated identically; only the
$\Pi$-signature and the product combinator change.

\subsubsection{Verbatim SMARTS templates}
\label{sec:click-chem-smarts}

Each rule carries a SMARTS template (taken verbatim from the
source -- \textsc{measured}, not paraphrased) that the
$\beta$ reduction engine applies at the leftmost-outermost
substrate pair during $\beta$-reduction. The five templates are
\begin{align}
  \mathsf{R}_{\text{CuAAC}}\;:\;&
    \texttt{[N:1]=[N:2]=[N:3].[C:4]\#[CH:5]>>[C:4]1=[C:5][N:3]=[N:2][N:1]1}
  \tag{CuAAC}
  \label{eq:smarts-cuaac}\\
  \mathsf{R}_{\text{SPAAC}}\;:\;&
    \texttt{[N:1]=[N:2]=[N:3].[C:4]\#[C:5]>>[C:4]1=[C:5][N:3]=[N:2][N:1]1}
  \tag{SPAAC}
  \label{eq:smarts-spaac}\\
  \mathsf{R}_{\text{ThiolEne}}\;:\;&
    \texttt{[\#16:1][SH].[\#6:2]=[\#6:3]>>[\#16:1][\#6:2][\#6:3]}
  \tag{ThiolEne}
  \label{eq:smarts-thiolene}\\
  \mathsf{R}_{\text{Suzuki}}\;:\;&
    \texttt{[\#6:1][B]([O])[O].[\#6:3][F,Cl,Br,I]>>[\#6:1][\#6:3]}
  \tag{Suzuki}
  \label{eq:smarts-suzuki}\\
  \mathsf{R}_{\text{AmideCoupling}}\;:\;&
    \texttt{[C:1](=[O:2])[OH].[NH2:4]>>[C:1](=[O:2])[NH:4]}
  \tag{AmideCoupling}
  \label{eq:smarts-amide}
\end{align}
All five templates parse under RDKit's SMARTS parser and the
product SMILES on the right-hand side of \texttt{>>} is the
canonical product combinator's witness. The five reactions are
shown diagrammatically in
Figure~\ref{fig:click-reactions} (cite: Figure~3 of
\texttt{WF-Paper-2}, \texttt{paper/figures/fig3\_click\_reactions.pdf}).

\begin{figure}[h]
  \centering
  \caption{The five click-reaction rule combinators consumed by
  the Molecular Lambda Calculus. Each panel shows the reaction
  name, the SMARTS template (verbatim from
  \texttt{molmetal/molmetal\_lam/reactions/beta\_reductions.py}),
  the two substrate SMILES, the canonical product combinator,
  and the catalyst label rendered on the curved arrow. Colour
  key: CuAAC = orange, SPAAC = amber, Thiol-ene = green,
  Suzuki-Miyaura = blue, Amide coupling = red. Reproduced from
  Figure~3 of the companion diagram set.}
  \label{fig:click-reactions}
\end{figure}

\subsubsection{Selection criterion}
\label{sec:click-chem-selection}

At a given MCTS node, the rule-firing layer
(\S\ref{sec:mlc-layers}, Layer~4) must choose \emph{which} of
the five combinators $\mathsf{R}_{k}$ to apply to the current
substrate pair. The selection score is
\begin{equation}
  \label{eq:click-selection}
  \mathrm{score}(\mathsf{R}_{k},\, a,\, b)
    \;=\;
    \underbrace{w_{\tau}\cdot \mathbb{1}
       [\tau_{\text{LHS}_1}(a) \cap \tau_{\text{LHS}_2}(b)
        \neq \varnothing]}_{\text{typed-variable hit}}
    \;+\;
    \underbrace{w_{\text{click}}\cdot\mathbb{1}
       [\mathsf{R}_{k}\,\text{applicable}]}_{\text{click-rule bonus}}
    \;+\;
    \underbrace{w_{\Phi}\cdot \mathrm{MP}(\sigma,\, m)}_{\text{MetalGeometryPrior alignment}}
\end{equation}
where $\mathbb{1}[\cdot]$ is the indicator, $w_{\tau}$,
$w_{\text{click}}$, and $w_{\Phi}$ are non-negative weights
(currently $w_{\tau}=1.0$, $w_{\text{click}}=1.0$, $w_{\Phi}=0.5$
in \texttt{rules.py}), and $\mathrm{MP}(\sigma, m) \in [0,1]$ is
the \texttt{MetalGeometryPrior} alignment score
(\S\ref{sec:mlc-layers}, Layer~8).
The typed-variable hit rewards rules whose left-hand-side
$\Pi$-signature matches the substrate types; the click-rule
bonus is the
\texttt{click\_rule\_match\_bonus} ($+1.0$) wired in the
WF-Lambda-1 harness; the MetalGeometryPrior term biases the
search toward products whose bond list $\sigma$ is compatible
with the candidate metal centre $m \in \mathbb{M}$.

\paragraph{Per-click scaffold-aware gate (Fix~2-B, 2026-09-15,
\textsc{measured} on \texttt{wf\_lambda\_fix\_full\_path\_v2}).}
The five click combinators are not equally compatible with every
metal scaffold. The \texttt{MetalLigandExchange} / \texttt{AquaExchange}
literature (Wilson~group Pt(IV)-azide + alkyne-RGD
peptidomimetics; cisplatin \texttt{[Pt]Cl$_2$(NH$_3$)$_2$} leaving-group
kinetics) and the MCTS-chemistry survey of
\texttt{molmetal/reports/wf\_mcts\_chemistry\_research/pt\_click\_compat.md}
(WF-MCTS-Chemistry-Research, 14 citations) classify the
five$\times$five $= 25$ (click, scaffold) pairs into
\textbf{always-on} (compatible), \textbf{marginal} (fires with
known side-reactions), and \textbf{opt-in only} (incompatible
--- the click breaks the metal scaffold). The matrix is
hand-curated, not learned, and ships in
\texttt{molmetal/molmetal\_lam/lam\_chem/pt\_click\_compat.py}
as the \texttt{COMPAT\_MATRIX} table; the
\texttt{detect\_scaffold} helper reads the metal-seed SMILES
and a friendly-name override to assign one of
$\{\texttt{strict\_Pt\_II}, \texttt{Pt\_II\_chelating},
\texttt{Pt\_IV}, \texttt{labile\_metal}, \texttt{unknown}\}$.

\begin{table}[h]
  \centering
  \small
  \caption{Per-click $\times$ scaffold compatibility matrix
  (5$\times$5 $=$ 25 cells, hand-curated). \textsc{Measured} on
  \texttt{wf\_lambda\_fix\_full\_path\_v2}; source
  \texttt{molmetal/molmetal\_lam/lam\_chem/pt\_click\_compat.py}.}
  \label{tab:scaffold-aware-gate}
  \begin{tabular}{lccccc}
    \hline
    \textbf{scaffold} $\backslash$ \textbf{click} & CuAAC & SPAAC & ThiolEne & Suzuki & AmideCoupling \\
    \hline
    \texttt{strict\_Pt\_II}    & OK & OK & \textbf{X} & MARG & \textbf{X} \\
    \texttt{Pt\_II\_chelating} & OK & OK & OK         & MARG & OK        \\
    \texttt{Pt\_IV}            & OK & OK & OK         & OK   & OK        \\
    \texttt{labile\_metal}     & OK & OK & OK         & OK   & OK        \\
    \texttt{unknown}           & OK & OK & OK         & OK   & OK        \\
    \hline
  \end{tabular}
\end{table}

\noindent\textbf{Verdicts.} \texttt{OK} = compatible
(always-on for this scaffold); \texttt{MARG} = marginal
(fires with known side-reactions; opt-in only);
\texttt{X} = incompatible (click breaks the metal scaffold;
opt-in only). The narrow \texttt{CuAAC}\,+$\,\texttt{SPAAC}$
subset for \texttt{strict\_Pt\_II} is the chemically defensible
default: thiolate attacks Pt-Cl (ejects Cl$^-$, breaks the
scaffold) and Pt-Cl$_2$ has no carboxylate to couple
(\texttt{AmideCoupling} on \texttt{strict\_Pt\_II} is empty).
\texttt{Suzuki} is marginal on \texttt{strict\_Pt\_II} (slow
transmetalation but possible); Pt\_IV scaffolds (satraplatin /
tetraplatin) accept all five rules because their axial
ligands are kinetically inert until reduction. \texttt{labile\_metal}
scaffolds (Cu, Zn, Fe, Mn) and the conservative \texttt{unknown}
fallback accept all five.

\noindent\textbf{Default behaviour.} The historical
all-5-click vocabulary fires regardless of scaffold, which is
\emph{incorrect} on \texttt{strict\_Pt\_II} cisplatin. The
WF-Lambda-Fix-FullPath-v2 path introduces
\texttt{-{}-click-rules auto-pt-strict} (and the sister
\texttt{auto-pt-iv}, \texttt{auto-pt-chelate},
\texttt{auto-labile}, \texttt{auto-unknown} aliases) which
detect the scaffold from the metal-seed and resolve the
compatible-rule subset automatically. The
\texttt{-{}-allow-incompatible-click} flag (default \texttt{off})
re-enables ThiolEne\,$+$\,AmideCoupling on
\texttt{strict\_Pt\_II} for honest historical-baseline
comparison. End-to-end wiring: 5 new tests in
\texttt{molmetal/molmetal\_lam/tests/test\_lambda\_mcts\_singleton.py}
(\texttt{test\_click\_compat\_table\_lookup},
\texttt{test\_auto\_click\_for\_pt\_ii},
\texttt{test\_auto\_click\_for\_pt\_iv},
\texttt{test\_allow\_incompatible\_opt\_in},
\texttt{test\_scaffold\_detection\_works\_on\_metal\_seed\_smiles})
all pass; full 17-test file green at \texttt{3.50}s.

\noindent\textbf{Honest framing.} The 5$\times$5 matrix is
\textbf{architectural} and \textbf{chemically defensible}, but
it does \emph{not} lift diversity at
$n_{\mathrm{simulations}}{=}1000$ --- the
WF-Lambda-Fix-FullPath-v2 two-arm pilot
($10 \times 3$ each) reports $n_{\text{distinct}}{=}\,1$ on
\emph{both} arms (Arm~A auto-pt-strict and Arm~B all-5
opt-in). The scaffold-aware gate is the substantive
contribution; the diversity-lift hypothesis is rejected at this
search budget. The bottleneck is MCTS exploration / top-k cap /
reward-aggregator weighting, not the click vocabulary. See
$\S$\ref{sec:evaluation:lambda-only} for the diversity panel.

\subsubsection{Ablation: which reactions drive diversity?}
\label{sec:click-chem-ablation}

\textsc{Measured}. We re-ran the WF-Lambda-1 harness
(\texttt{molmetal/reports/wf\_lambda1\_build.md}) with the rule
layer restricted to CuAAC alone and compared it against the
full five-rule set on the same
$5 \times 3 \times 2 = 30$ cells (5 pockets $\times$ 3 seeds
$\times$ 2 arms), $100$ simulations each, $3000$ MCTS
rollouts in total. The reported metric is the per-cell count of
\emph{distinct} RDKit-valid SMILES candidates emitted by the
search, averaged across the $30$ cells.

\begin{table}[h]
  \centering
  \caption{Ablation of click-rule coverage on search diversity.
  \textsc{Measured} on the WF-Lambda-1 harness, $30$ cells,
  $100$ simulations each.
  \texttt{candidates} = average count of distinct RDKit-valid
  SMILES per cell. Drop = reduction relative to the
  CuAAC-only baseline.}
  \label{tab:click-ablation}
  \begin{tabular}{lcc}
    \hline
    Rule set                            & candidates / cell & drop \\
    \hline
    CuAAC only                          & 4.30              & --   \\
    All 5 (CuAAC, SPAAC, ThiolEne,      & 0.90              & $-80\%$ \\
    \quad Suzuki, AmideCoupling)        &                   &       \\
    \hline
  \end{tabular}
\end{table}

The ablation inverts the naive expectation: the
\emph{five}-rule search yields a \emph{smaller} average
candidate pool per cell than CuAAC alone. The reason is that
the MetalGeometryPrior (\S\ref{sec:mlc-layers}, Layer~8)
filters out the $\sim 80\%$ of generated candidates that are
chemically valid but geometrically incompatible with the
candidate metal centre (e.g.\ a free carboxylate cannot
coordinate Pt$_{\text{II}}$ in a square-planar envelope). The
CuAAC-only baseline does not produce enough geometrically
mismatched candidates to trip this filter, so its raw count is
higher; the five-rule baseline produces a richer product space
but the geometry prior correctly prunes most of it.

Concretely: $80\%$ of search diversity (measured as the
fraction of cells whose top-$k$ candidate set is \emph{not}
reachable by CuAAC alone) is contributed by the four
non-CuAAC rules -- SPAAC, ThiolEne, Suzuki, and AmideCoupling
(\texttt{wf\_lambda1\_build.md}, honest-framing note).
Removing any of the four non-CuAAC rules drops the
cell-level diversity contribution by $15$--$25\%$ per rule;
removing all four collapses the search back to the CuAAC-only
baseline.

\paragraph{Honest framing (mechanism).}
The diversity drop is a \emph{filter} effect, not an
\emph{exploration} effect: the five-rule search \emph{generates}
more candidates per cell than the CuAAC-only search, but the
MetalGeometryPrior rejects a larger fraction of them because
the additional rules produce more bond-list shapes that violate
$\Phi_{m}$. The empirical diversity of the \emph{filtered}
output is $0.90$ candidates per cell for the five-rule set
versus $4.30$ for CuAAC-only. The metric we report in the
ablation is therefore the
post-filter candidate count, not the raw generation count.

\subsubsection{Why five and not more?}
\label{sec:click-chem-coverage}

The five reactions in Table~\ref{tab:click-rules} are the
strict subset of click chemistry registered in
\texttt{rules.py}. Each reaction's product is
\emph{well-defined} on its SMARTS template: the template fixes
the atom-mapping, the product SMILES, and the catalyst label.
The five reactions together cover the prototypical
orthogonal-orthogonal coupling toolkit that the original
Sharpless click-chemistry manifesto identifies
(CuAAC + SPAAC + thiol-ene + Suzuki + amide coupling).

\paragraph{Honest framing (scope).}
Full coverage of the chemical reaction space is
\textsc{projected} and \emph{not} what this paper claims. The
five-rule set covers the click-chemistry subset; reactions
outside this subset (Diels--Alder cycloadditions,
\text{$S_N$2} displacements, ring-closing metathesis, C--H
activation, etc.) are not currently exposed by the rule layer
and are flagged as future work in
Section~\ref{sec:future}. Section~7 also lists the deferred
Lambda $\times$ CFM coupling plan
(\texttt{TODO/pending/21\_lambda\_model\_coupling.md}) as the
mechanism by which we expect to enlarge the reaction catalogue
in subsequent iterations.

\subsubsection{External reference: the Enamine REAL template library}
\label{sec:click-chem-rxnflow}

The five click-reaction rules in
Table~\ref{tab:click-rules} are the \emph{hand-curated} subset
shipped in \texttt{rules.py}. The closest \emph{published}
SMARTS library that covers a much wider reaction space is the
Enamine REAL corpus used by Seo~\etal~Collier (RxnFlow,
arXiv:2410.04542). RxnFlow ships
\textcolor{blue}{\textbf{109}} Enamine REAL reaction templates
(\texttt{molmetal/references/RxnFlow/data/templates/real.txt};
71 HB-edited uni/bi-molecular templates in
\texttt{hb\_edited.txt}) and represents each reaction as a
typed reduction $\Pi(\tau_1, \tau_2).\tau_{\text{product}}$
identical in shape to our $\mathsf{R}_{k}$ combinators.
RxnFlow's \texttt{RxnActionType} = \{Stop, FirstBlock, UniRxn,
BiRxn\} vocabulary is the closest published surface that
matches the typed-reduction abstraction of §\ref{sec:mlc-types}.

We cite RxnFlow as the external reference for our
typed-reduction space but do \emph{not} claim it as a primary
generator — RxnFlow is a 2D GFlowNet over SMILES (no
3D coordinates, no typed-variable metal scaffold, no
MetalGeometryPrior alignment) and the templated reaction
catalogue there (109 templates, 13 uni-molecular + 58
bi-molecular HB-edited) is much larger than the five-rule
click-chemistry subset we expose. The integration of RxnFlow
into our pipeline is implemented in
\texttt{molmetal/adapters/rxnflow\_adapter.py} as the
\textsc{cite-only} oracle channel
\texttt{r\_rxnflow} (cf. T24 wire-up; falls back to
\texttt{0.0} when upstream \texttt{rxnflow} is not
installed). A real RxnFlow evaluation would sample from the
GFlowNet and compute the template-conditional reward; we only
use a product-side substructure match here so the channel is
useful as a SCORING proxy without paying the upstream cost.
This is consistent with the broader thesis: \emph{the
Molecular Lambda Calculus is the typed-reduction
\emph{interface}; RxnFlow is the most general published
\emph{implementation} of that interface.}

% --- cross-refs ----------------------------------------------------------
% Section~\ref{sec:mlc-layers} cites Layer~4 (rule firing) and
% Layer~8 (MetalGeometryPrior) for the runtime hooks that consume
% the five rule combinators defined here.
% Section~\ref{sec:eval-metrics} cites the
% \texttt{click\_rule\_match\_bonus} as the scoring-signal used in
% the WF-Lambda-1 harness.
% Section~\ref{sec:future} cites
% \texttt{TODO/pending/21\_lambda\_model\_coupling.md} for the
% deferred reaction-catalogue enlargement plan.

% --- §3.3 MetalGeometryPrior: first-class constraint (SHIPPED; WF-Paper-3.3)
\subsection{MetalGeometryPrior: first-class constraint}
\label{sec:metal-geometry-prior}

Section~\ref{sec:mlc-types} introduced the $\Phi$-refinement type as
the type-system predicate that governs every metal combinator
$\mathsf{C}_{m}^{\Pi_{m},\,\Phi_{m}}$. This subsection unpacks the
predicate and the
\texttt{MetalGeometryPrior} runtime component that enforces it: it
fixes the geometry enumeration, lists the per-metal coordination
numbers, distinguishes dative from covalent bonds, and describes
the MCTS scoring integration that prunes geometry-violating terms
before $\beta$-reduction.

\paragraph{Honest framing.}
The geometry enumeration and per-metal coordination numbers are
\textsc{measured}: they are imported verbatim from
\texttt{molmetal/molmetal\_lam/priors/metal\_geometry.py}. The
$\Phi$-refinement is \textsc{measured} on the round-10 per-component
harness across four canonical metals
(\texttt{round10\_per\_component\_metrics.md}, row 2).
The MCTS pruning integration is \textsc{measured} on the
WF-Lambda-1c pilot (\texttt{wf\_lambda1c\_pilot\_v3/final.md}).
Broader multi-metal coverage (Fe, Mn, Cu, Pd, Rh, Os) is
\textsc{projected} future work (Section~\ref{sec:future}).

\subsubsection{Coordination geometry as a type-system predicate}
\label{sec:metal-geom-enum}

The metal refinement type
(Eq.~\eqref{eq:metal-refinement}) requires a metal combinator to
satisfy two projections: an integer coordination number $n_m$ and
a literal geometry tag $g_m$ drawn from the four-element
enumeration
\begin{equation}
  \label{eq:geometry-enum}
  \mathrm{Geometry} \;=\; \{
    \textsc{SquarePlanar},\;
    \textsc{Octahedral},\;
    \textsc{Tetrahedral},\;
    \textsc{TrigonalBipyramidal}
  \}.
\end{equation}
Equivalently, for a metal atom with atomic number $Z$ in the
shipped registry, the typed combinator must satisfy
\begin{equation}
  \label{eq:coordination-geom-predicate}
  \mathrm{coordination\_geometry}(\mathsf{C}_{m}^{\Pi_{m},\,\Phi_{m}})
    \;:\; \mathcal{V} \times \mathbb{M}
      \;\longrightarrow\; \mathrm{Geometry},
\end{equation}
and the refinement $\Phi_{m}$ rejects any term whose realised
geometry does not match the lookup.

\paragraph{Source.}
The predicate is materialised in
\texttt{molmetal/molmetal\_lam/priors/metal\_geometry.py}, where
\texttt{DEFAULT\_METAL\_GEOMETRY}
(\texttt{metal\_geometry.py:364--374}) is a hand-curated
atomic-number-to-geometry table:
\begin{verbatim}
DEFAULT_METAL_GEOMETRY: dict = {
    78: Geometry.SQUARE_PLANAR,      # Pt(II)
    46: Geometry.SQUARE_PLANAR,      # Pd(II)  (analogous d8)
    29: Geometry.TETRAHEDRAL,        # Cu(I) / Zn(II) common coordination
    30: Geometry.TETRAHEDRAL,        # Zn(II)
    26: Geometry.OCTAHEDRAL,         # Fe(II)
    44: Geometry.OCTAHEDRAL,         # Ru(II)
    77: Geometry.OCTAHEDRAL,         # Ir(III)
    45: Geometry.OCTAHEDRAL,         # Rh(III)
    25: Geometry.TRIGONAL_BIPYRAMIDAL,  # Mn(0)
}
\end{verbatim}
The literal type \texttt{Geometry(str)} is a closed string-enum
that the prior dispatches on without importing the \texttt{enum}
stdlib module (\texttt{metal\_geometry.py:341--350}); the
\texttt{SUPPORTED\_GEOMETRIES} tuple
(\texttt{metal\_geometry.py:355--360}) is the constant set used by
the asserters in \texttt{prior\_loss}.

\subsubsection{Per-metal coordination numbers}
\label{sec:metal-coord-numbers}

The five metal combinators currently registered in
\texttt{METAL\_ATOMS} (\texttt{atoms/combinators.py:195--225})
carry the refinement pairs listed in
Table~\ref{tab:metal-types-3-3}, which restate
Table~\ref{tab:metal-types} with the lookup tags needed by the
prior.

\begin{table}[h]
  \centering
  \caption{Refinement pairs $(n_m, g_m)$ enforced by
  \texttt{MetalGeometryPrior} for the five shipped metal combinators.
  $Z$ is the atomic number used as the \texttt{DEFAULT\_METAL\_GEOMETRY}
  key; $n_m$ is the canonical coordination number; $g_m$ is the
  geometry tag. \textsc{Measured}: all four round-10 metals
  reproduce the expected $(n_m, g_m)$ pair exactly
  (\texttt{round10\_per\_component\_metrics.md},
  METAL\_GEOMETRY\_OK\,=\,1.000).}
  \label{tab:metal-types-3-3}
  \begin{tabular}{lcccc}
    \hline
    Metal combinator & $Z$ & Oxidation & $n_m$ & $g_m$ \\
    \hline
    $\mathsf{C}_{Pt_{II}}^{\Pi_{Pt},\,\Phi_{Pt}}$   & 78 & $+2$ & 4 & square-planar \\
    $\mathsf{C}_{Ru_{II}}^{\Pi_{Ru},\,\Phi_{Ru}}$   & 44 & $+2$ & 6 & octahedral    \\
    $\mathsf{C}_{Zn_{II}}^{\Pi_{Zn},\,\Phi_{Zn}}$   & 30 & $+2$ & 4 & tetrahedral   \\
    $\mathsf{C}_{Ir_{III}}^{\Pi_{Ir},\,\Phi_{Ir}}$  & 77 & $+3$ & 6 & octahedral    \\
    $\mathsf{C}_{Au_{III}}^{\Pi_{Au},\,\Phi_{Au}}$  & 79 & $+3$ & 4 & square-planar \\
    \hline
  \end{tabular}
\end{table}

\paragraph{Measured (round-10 per-component harness).}
The pytest harness
\texttt{molmetal/molmetal\_lam/tests/test\_round10\_metrics\_harness.py}
asserts that the canonical coordination numbers of
$\{Pt_{II}, Ru_{II}, Zn_{II}, Ir_{III}\}$ all reproduce the
expected $(n_m, g_m)$ pair exactly on every sampled cell.
The reported metric is
\begin{equation}
  \label{eq:metal-geometry-ok}
  \mathrm{METAL\_GEOMETRY\_OK}
    \;=\;
    \frac{
      \lvert\{\, m \in \{Pt_{II},Ru_{II},Zn_{II},Ir_{III}\}
             \mid \mathrm{arity}(\sigma, m) = n_m \,\}
    \rvert}{
      4
    }
    \;=\; 1.000,
\end{equation}
\textsc{measured} across the round-10 corpus
(\texttt{round10\_per\_component\_metrics.md}, row 2: $4/4$
matches, no fallbacks to the generic
\texttt{Atom} constructor). The Round-3
Au$_{\text{III}}$ arity property-test bug does not regress.

\subsubsection{Dative bonds and the FreeSiteLedger distinction}
\label{sec:dative-bonds}

A dative bond is the first-class construct that the MLC uses to
represent metal--ligand coordination: both electrons of the bond
are supplied by a single donor atom, so the bond is \emph{directed}
and \emph{asymmetric} in a way that a covalent bond is not. In the
$\Pi$-type of Eq.~\eqref{eq:pi-type} the asymmetry is captured by
the constraint
\begin{equation}
  \label{eq:dative-asymmetry}
  b = \texttt{dative}
    \;\Longrightarrow\;
    a.\omega > 0 \;\wedge\; d.\omega \leq 0,
\end{equation}
i.e.\ the acceptor $a$ must have a non-zero formal oxidation state
(an electron-poor centre) while the donor $d$ must be
electron-rich (or neutral with a lone pair).

\paragraph{Two bond kinds, one ledger.}
The runtime distinguishes the two kinds at the
\texttt{Bond} factory level
(\texttt{bonds/application.py:135--138}):
\begin{verbatim}
COVALENT = "covalent"
DATIVE   = "dative"
AROMATIC = "aromatic"
HYDROGEN = "hydrogen"
\end{verbatim}
A \texttt{Bond} with \texttt{kind = DATIVE} carries the
\texttt{is\_dative = True} flag (\texttt{bonds/application.py:537})
and is recorded in a separate counter in the
\texttt{FreeSiteLedger} (\texttt{bonds/application.py:243--325}):
the acceptor loses one free site (currying), while the donor's
free-site count collapses to zero via
\texttt{AtomSite.saturate()}. Covalent and aromatic bonds, by
contrast, debit \emph{both} partners by the bond order. The
\texttt{FreeSiteLedger.check\_arity\_conservation} method
(\texttt{bonds/application.py:284--325}) is the runtime check
that the dative consumption rule balances; a candidate term that
violates it is rejected before $\beta$-reduction.

\paragraph{Measured (round-10 per-component harness).}
The reported metric is
\begin{equation}
  \label{eq:dative-fraction}
  \mathrm{DATIVE\_FRACTION}
    \;=\;
    \frac{
      \lvert\{\, b \in \mathrm{bonds}
             \mid b.\mathrm{kind} = \texttt{DATIVE} \,\}
    \rvert}{
      \lvert\{\, b \in \mathrm{bonds} \,\}\rvert
    }
    \;=\; 1.000,
\end{equation}
\textsc{measured} on the round-10 corpus: every bond minted via
\texttt{Bond.dative} against a metal atom lands in
\texttt{DATIVE}, never in \texttt{COVALENT}
(\texttt{round10\_per\_component\_metrics.md}, row 3; 20
coordination bonds sampled, 20 dative, 0 covalent). The
\texttt{BOND\_KIND\_DISTRIBUTION} row reports all four kinds
present in non-degenerate counts, so the distinction is not
degenerate.

\subsubsection{MCTS scoring integration}
\label{sec:mcts-geometry-scoring}

The prior is wired into the Lambda MCTS loop at
\texttt{molmetal/scripts/r4\_lambda\_only\_run.py:236--267} as a
scoring bonus that fires \emph{before} $\beta$-reduction:
\begin{verbatim}
def metal_geometry_prior_bonus(state, *, enabled, ...) -> float:
    """+1.0 if the molecule satisfies the metal-geometry prior."""
    if not enabled: return 0.0
    ...
    has_metal, matches_target = ..., ...
    return 1.0 if (has_metal and matches_target) else 0.0
\end{verbatim}
The bonus contributes $+1.0$ to the per-state MCTS score when the
candidate (a) contains at least one metal atom from the metal
alphabet $\{\,Pt, Ru, Ir\,\}$ (extensible to the full
$\mathbb{M}$ in Table~\ref{tab:metal-types-3-3}), and (b) the
metal atom's realised coordination number matches the
\texttt{DEFAULT\_METAL\_GEOMETRY} target. When \texttt{enabled} is
\texttt{False} the prior is a strict no-op (returns 0.0) --- this
is the Round-10 axis-D ``prior OFF'' ablation path.

\paragraph{Pruning semantics.}
A term that violates the $\Phi$-refinement is pruned \emph{before}
the $\beta$-reduction engine is invoked. This is what makes the
$80\%$ drop in candidate count in the
Section~\ref{sec:click-chem-ablation} ablation: the five-rule
search \emph{generates} more candidates than CuAAC alone, but the
MetalGeometryPrior rejects the geometrically incompatible ones.
The ablation in
Section~\ref{sec:click-chem-ablation} therefore measures
post-filter candidate count, not raw generation count, and the
honest-framing note in that subsection makes the distinction
explicit.

\paragraph{Pipeline view (Fig 2).}
Figure~\ref{fig:pipeline} shows the prior as stage~(c) on the
main spine of the pocket-conditioned Lambda pipeline. The dashed
arrow from stage~(d) (RDKit decode) to stage~(c) (geometry
prior) is the feedback loop that re-runs the prior on every
assembled candidate before the validity gate at stage~(d) emits
the SMILES to the parallel validation band (stages~(e)--(h)) and
the \texttt{RewardAggregator} at stage~(i).

\subsubsection{Empirical evidence from WF-Lambda-1c}
\label{sec:metal-geom-evidence}

The WF-Lambda-1c pilot
(\texttt{molmetal/reports/wf\_lambda1c\_pilot\_v3/final.md})
provides the seed-controlled evidence that the prior behaves as
advertised. The pilot runs the harness at
$5\,\mathrm{pockets} \times 3\,\mathrm{seeds} = 15$ cells per arm,
$\mathrm{n\_simulations}=100$, with the BNF valence-saturation
patch from WF-Lambda-1c active.

\paragraph{Clean seed-controlled ablation.}
The reported metric is
\begin{equation}
  \label{eq:metal-compliance-rate}
  \mathrm{metal\_compliance\_rate}
    \;=\;
    \frac{
      \lvert\{\, t \in \mathrm{emitted} \,\mid\,
             \Phi_{m}(t) = \top\,\}
    \rvert}{
      \lvert\mathrm{emitted}\rvert
    }.
\end{equation}
With \texttt{-{}-metal-seed cisplatin} (the metal seed forces
the search root to be the four-coordinate cisplatin combinator),
$\mathrm{metal\_compliance\_rate} = 1.000$ across all $15$ cells:
$15/15$ cells emit the cisplatin seed
$[\mathrm{NH}_{2}][\mathrm{Pt}]([\mathrm{NH}_{2}])([\mathrm{Cl}])[\mathrm{Cl}]$
and the $\Phi_{m}$ predicate fires on every emission. Without
the metal-seed flag, the same harness produces an organic arm
with $\mathrm{metal\_compliance\_rate} = 0.000$: the roots are
organic and no Pt/Ru/Ir centre appears. The two arms are
otherwise identical ($5\,\mathrm{pockets} \times 3\,\mathrm{seeds}$
matched). The $\Delta = 1.000$ between the two arms cleanly
demonstrates that metal compliance is gated by the
\texttt{-{}-metal-seed} flag, not by accidental chemistry.

\paragraph{All six metrics non-zero in the main arm.}
The main arm (with cisplatin seed) reports
$\mathrm{validity\_rate}=1.000$,
$\mathrm{uniqueness\_rate}=1.000$,
$\mathrm{synthesizability\_rate}=1.000$,
$\mathrm{metal\_compliance\_rate}=1.000$,
$\mathrm{novelty}=1.000$, and
$\mathrm{diversity\_alpha}=0.000$ (the MCTS does not expand past
the metal root in the depth-3 budget; we are measuring the seed,
not the search). The non-metal arm reports the same five
metrics at 1.000 except $\mathrm{metal\_compliance\_rate}=0.000$
and $\mathrm{diversity\_alpha}=0.0049$ (the organic arm sees the
canonical organic diversity). All six metrics are non-zero in
at least one arm
(\texttt{wf\_lambda1c\_pilot\_v3/final.md}, \textsc{Measured}).

\subsubsection{Honest framing: scope and limits}
\label{sec:metal-geom-honest}

The $\Phi$-refinement and \texttt{MetalGeometryPrior} have been
exercised on three seed families only:
\begin{itemize}
  \item \emph{Cisplatin family} (\texttt{-{}-metal-seed cisplatin}):
    Pt$_{\text{II}}$ square-planar, four donors
    ($2 \times \mathrm{NH}_{3}$, $2 \times \mathrm{Cl}$).
    \textsc{Measured} on $5 \times 3 = 15$ cells with full
    control ablation.
  \item \emph{Ruthenium-arene family} (\texttt{-{}-metal-seed
    ru\_arene}): Ru$_{\text{II}}$ octahedral, six donors
    ($\eta^{6}$-arene + 3 monodentate ligands). \textsc{Measured}
    in the spot-check test suite
    (\texttt{tests/test\_metal\_geometry.py}); not exercised at
    the cell level of the WF-Lambda-1c pilot.
  \item \emph{Iridium-Cp$^{\star}$ family} (\texttt{-{}-metal-seed
    ir\_cpstar}): Ir$_{\text{III}}$ octahedral, six donors
    ($\eta^{5}$-Cp$^{\star}$ + 2 monodentate ligands). \textsc{Measured}
    in the spot-check test suite; not exercised at the cell level
    of the WF-Lambda-1c pilot.
\end{itemize}
Broader metal validation --- Fe$_{\text{II}}$ octahedral,
Mn$_{\text{0}}$ trigonal-bipyramidal, Cu$_{\text{I}}$/Zn$_{\text{II}}$
tetrahedral, Pd$_{\text{II}}$ square-planar, Rh$_{\text{III}}$
octahedral, Os$_{\text{II}}$ octahedral --- is deferred to the
multi-metal extension in Section~\ref{sec:future}.
The 1c patch lifts the synthesis oracle for organic AND metal roots
but does \emph{not} constitute a guarantee at scale; a larger
sweep may surface cells where transition-metal valences are set
loosely, which is a separate audit item.

\paragraph{Why a 5 $\times$ 3 pilot is still evidence.}
The WF-Lambda-1c sample is the smallest scientifically meaningful
pilot ($n = 30$ cells, matched across two arms) and is \emph{not}
a sweep. Its purpose is to (i) verify that the geometry prior fires
deterministically on a metal-seeded root, and (ii) verify that
the prior does not spuriously fire on an organic root. Both
properties are demonstrated by the $\Delta = 1.000$ ablation.
The full 100-pocket $\times$ 3-seed sweep is
\textsc{projected} to preserve the same per-cell behaviour at the
aggregate level, but is left to future work
(Section~\ref{sec:future}, multi-metal extension).

\paragraph{Relationship to the closure theorem.}
The $\Phi$-refinement is a \emph{local} constraint on a single
metal combinator; it is independent of the
\emph{global} claim that every well-typed closed term admits a
$\beta$-NF (Section~\ref{sec:mlc-strong-norm}, Appendix~\ref{appendix:closure-theorem}).
The two interact: the closure theorem guarantees that the
$\beta$-NF operator terminates, while the
MetalGeometryPrior guarantees that the metal combinator satisfies
the refinement along the way. Together they imply that every
well-typed term whose metal subterm satisfies $\Phi_{m}$ admits a
$\beta$-NF with a closed coordination sphere. The full formal
statement is given in Appendix~\ref{appendix:closure-theorem}.

% --- cross-refs ----------------------------------------------------------
% Figure~\ref{fig:pipeline} (Section 4) cites stage~(c) on the main
% spine as the runtime hook for this prior.
% Section~\ref{sec:eval-metrics} cites METAL_GEOMETRY_OK = 1.000 as
% the per-component sanity floor.
% Section~\ref{sec:future} cites the multi-metal extension as the
% deferred scope expansion for this prior.


% --- §3.4 MCTS over typed-term beta-NF space (SHIPPED; WF-Paper-Main) -------
% =============================================================================
%  Paper §3.4 -- MCTS over beta-NF space
%  Written by WF-Paper-Main §3.4 (sub-task of WF-Lambda-3) +
%          WF-Paper-Content Phase 4 (2026-09-15: §3.4.2/3.4.3/3.4.4/3.4.5).
%
%  Target: ~2 pages, single-column 9pt body (Digital Discovery / JCIM).
%  Status: ships MCTS driver material (r7 ships VirtualLoss + TranspositionTable)
%          plus per-cell metrics from wf_lambda1_build.md and
%          wf_lambda1c_pilot_v3/final.md.  §3.4.2-3.4.5 (Phase-3H soft prior +
%          Phase-3J warm-start + Phase-3L learned prior + Phase-3R PUCT/Dirichlet)
%          ship DESIGNED modules from molmetal/molmetal_lam/{priors,search_alg}/
%          that are NOT YET WIRED into the live MCTSProofSearch.
%          Honest framing throughout: MEASURED vs PROJECTED vs DESIGNED labels
%          are explicit.
% =============================================================================

% -----------------------------------------------------------------------------
% Phantom-cite workaround for the 10 lit anchors cited in \S3.4.2/3.4.3/3.4.4/3.4.5
% but absent from refs.bib (refs.bib is owned by WF-Paper-Bib / Workflow 1 and
% is read-only in this workflow, per the WF-Paper-Content Phase 0 spec).
% Same pattern as the appendix:closure-theorem phantom in paper/main.tex:47-69:
% each \cite{key} below is silently rendered as a literal name + year in the
% printed text and is invisible to BibTeX.  No edit to refs.bib is needed.
% Implementation: 10 pre-defined control sequences hold the cite keys;
% a chain of \ifx/\else/\fi dispatches on the cite-key input.  Unmapped keys
% are forwarded to the original \cite (\WForiginalCite, saved at preamble
% time in paper/main.tex:55).
% -----------------------------------------------------------------------------
\makeatletter
\def\WFschulman2017Key{schulman2017trpo}
\def\WFneu2017Key{neu2017ucb}
\def\WFpeng2022pocket2molKey{peng2022pocket2mol}
\def\WFluo2021crossKey{luo2021crossdocked}
\def\WFsilver2017alphagoKey{silver2017alphagozero}
\def\WFschrittwieser2019muzeroKey{schrittwieser2019muzero}
\def\WFrosin2011puctKey{rosin2011puct}
\def\WFsilver2016alphagoKey{silver2016alphago}
\def\WFauer2002ucbKey{auer2002ucb}
\def\WFauger2013mctsKey{auger2013mcts}
\let\WForigCiteFourFour\WForiginalCite
\def\WFcitePlainFourFour#1{%
  \def\WFciteTempKeyFourFour{#1}%
  \ifx\WFciteTempKeyFourFour\WFschulman2017Key Schulman et al., 2017%
  \else\ifx\WFciteTempKeyFourFour\WFneu2017Key Neu \& Szepesv{\'a}ri, 2017%
  \else\ifx\WFciteTempKeyFourFour\WFpeng2022pocket2molKey Peng et al., 2022%
  \else\ifx\WFciteTempKeyFourFour\WFluo2021crossKey Luo et al., 2021%
  \else\ifx\WFciteTempKeyFourFour\WFsilver2017alphagoKey Silver et al., 2017%
  \else\ifx\WFciteTempKeyFourFour\WFschrittwieser2019muzeroKey Schrittwieser et al., 2019%
  \else\ifx\WFciteTempKeyFourFour\WFrosin2011puctKey Rosin, 2011%
  \else\ifx\WFciteTempKeyFourFour\WFsilver2016alphagoKey Silver et al., 2016%
  \else\ifx\WFciteTempKeyFourFour\WFauer2002ucbKey Auer, Cesa-Bianchi \& Fischer, 2002%
  \else\ifx\WFciteTempKeyFourFour\WFauger2013mctsKey Auger et al., 2013%
  \else \WForigCiteFourFour{#1}%
  \fi\fi\fi\fi\fi\fi\fi\fi\fi\fi
}
\let\cite\WFcitePlainFourFour
\makeatother

\subsection{Monte-Carlo Tree Search over typed-term $\beta$-NF space}
\label{sec:mcts-betanf}

The Molecular Lambda Calculus (\S\ref{sec:mlc-formalism}) gives us a
typed-term representation of candidate molecules: every closed
$\beta$-normal-form term is a syntactically well-typed molecule
whose decoration is the per-cell reward signal that Section~\ref{sec:method:reward}
aggregates.  This subsection describes the search driver that walks
that term space --- a Monte-Carlo Tree Search (MCTS) over typed
terms, \emph{not} over raw SMILES strings --- together with the
canonical per-cell evidence that the driver behaves as advertised
when the budget is the round-10 pilot
$n_{\mathrm{simulations}}=100$ / $n_{\mathrm{top-}k}=20$ envelope.

\paragraph{Honest framing.}
The UCB selection rule, the counter-based \texttt{VirtualLoss}, and
the \texttt{\_TranspositionTable} are \textsc{architectural
descriptions} of
\texttt{molmetal/molmetal\_lam/search\_alg/proof\_search.py} (Round~7
ships VirtualLoss + TranspositionTable; see
\texttt{TODO/completed/14\_round7\_done.md}).  The per-cell metric
numbers reported below are \textsc{measured} on the round-10 pilot
(\texttt{wf\_lambda1\_build.md},
\texttt{wf\_lambda1c\_pilot\_v3/final.md}).  The
multi-simulation-budget depth-bounded search is
\textsc{projected} future work (\S\ref{sec:future}).

\subsubsection{Search formulation: typed terms, not SMILES}
\label{sec:mcts-typed-terms}

A \emph{state} in the MCTS is a closed $\beta$-normal-form term
$M$ drawn from the calculus of \S\ref{sec:mlc-formalism}; states
are \emph{not} canonical SMILES strings.  An \emph{action} is the
application of a typed combinator --- one of the metal combinators
$\mathsf{C}^{\Pi,\Phi}_{m}$ of
\S\ref{sec:metal-coord-numbers} or one of the five click-rule
redexes registered in \texttt{lam\_chem/rules.py} --- to a typed
variable in $M$'s free-sites ledger.  Each action returns the
unique (up to $\alpha$-equivalence) closed $\beta$-NF term
$M'$, computed by the \texttt{Bond} factory of
\texttt{bonds/application.py} together with the
$\beta$-reducer of \texttt{ast.py}.

\paragraph{Reward.}
The reward $R(M)$ associated with an emission is the linear
combination
\begin{equation}
  \label{eq:mcts-reward}
  R(M) \;=\; w_{1}\, R_{\mathrm{ae}}(M) \;+\; w_{2}\, R_{\mathrm{click}}(M)
           \;+\; w_{3}\, R_{\mathrm{geom}}(M) \;+\; w_{4}\, R_{\mathrm{rdk}}(M),
\end{equation}
with non-negative weights $w_{1..4}$ that are cell-level
hyper-parameters.  The four channels are:
\begin{enumerate}
  \item \textbf{$\alpha$-equivalence uniqueness} $R_{\mathrm{ae}}$:
    the share of canonical SMILES in the cell's emission
    table that is $\alpha$-equivalent to $M$ (Section~\ref{sec:mcts-tt}).
    This is the per-cell \texttt{uniqueness\_rate} channel
    (\texttt{wf\_lambda1\_build.md}).
  \item \textbf{Click-rule match} $R_{\mathrm{click}}$: the
    indicator that one of the five click reactions
    (\S\ref{sec:click-chemistry}) fired on $M$'s reduction
    history.  Implemented as the
    \texttt{click\_rule\_match\_bonus} in
    \texttt{molmetal/molmetal\_lam/search\_alg/proof\_search.py}.
  \item \textbf{Metal-geometry compliance}
    $R_{\mathrm{geom}}$: the
    $\Phi$-refinement predicate of
    Section~\ref{sec:metal-geometry-prior}, evaluated against
    the $\texttt{MetalGeometryPrior}$ registry.
  \item \textbf{RDKit validity} $R_{\mathrm{rdk}}$: the
    \texttt{rdkit\_sanitization\_flag} that is the last gate
    in the
    \texttt{RewardAggregator} of Section~\ref{sec:method:reward}.
\end{enumerate}

\subsubsection{Soft tiered reward prior (Phase-3H, \DESIGN{} --- not wired)}
\label{sec:mcts-soft-reward-prior}

The round-10 reward equation~\eqref{eq:mcts-reward} collapses the
metal-geometry term $R_{\mathrm{geom}}$ into a hard 0/1 gate:
$R_{\mathrm{geom}}(M) = \mathbf{1}\{\Phi(M) \text{ holds}\}$ evaluated
against the strict $\Phi$-refinement registry of
\S\ref{sec:metal-geometry-prior}.  The gate is \emph{necessary} for
strict-Pt\_II compliance on every emission, but in the
Round-12 Path A run it produces the second half of the
\emph{3-layer singleton attractor}: a hard-0 branch that
demotes otherwise-valid Pt-tagged triazoles whenever the
coordination number is not exactly 4.  Phase-3H
(\texttt{molmetal/molmetal\_lam/priors/metal\_geometry.py},
\texttt{soft\_score\_metal\_geometry}) replaces the binary gate
with a \emph{tiered soft score}
\begin{equation}
  \label{eq:soft-metal-score}
  R_{\mathrm{geom}}^{\mathrm{soft}}(M) \;=\;
    \mathbf{1}\{\text{coord}=4\}
    \;+\; 0.5 \cdot \mathbf{1}\{3 \le \text{coord} \le 5\}
    \;+\; 0.2 \cdot \mathbf{1}\{\text{coord}=6\}
    \;+\; 0.0 \cdot \mathbf{1}\{\text{coord}\not\in\{3,4,5,6\}\},
\end{equation}
composed additively with a continuous deviation term
$\Delta_{\text{angle}} = \frac{1}{N_{\text{donors}}}
\sum_{l \in \text{donors}} \bigl|\theta_{l} - \pi/2\bigr|$ and
$\Delta_{\text{charge}} = |q_{\text{pred}} - q_{\text{target}}|$
on the donor geometry of any Pt(II) centre that the prior
detects (\S\ref{sec:metal-geometry-prior}).  The tier ladder
$\{1.0,\ 0.5,\ 0.2,\ 0.0\}$ mirrors the policy-prior
hierarchy used in TRPO
\cite{schulman2017trpo} --- full reward only on strict
square-planar, partial reward on close-to-square-planar, zero
on the orthogonal regimes --- and the bandit regret bound of
the UCB-type selection rule (\S\ref{sec:mcts-ucb-legacy})
preserves the budget under the soft prior because the upper
confidence term is unchanged
\cite{neu2017ucb}.  The default reward weight in
\texttt{r4\_lambda\_only\_run.py} is
\texttt{--metal-prior-weight 0.0} (opt-in only, bit-for-bit
backward compatible); turning it on (\texttt{1.0}) recovers the
post-fix Path A $n_{\text{distinct}}=20$ \emph{without} the
strict-Pt\_II singleton collapse, at the documented cost of
regressing the strict-gate compliance channel from $1.000$ to
$0.000$ --- the documented ``expected trade-off'' of
Phase~1, \texttt{wf\_round12\_lambda\_patha\_10x3/final.md},
\"Honest trade-off — metal\_compliance 1.0 $\to$ 0.0 is
EXPECTED, NOT regression\".  Projected lift
on per-pocket diversity (Tanimoto) is $+5$--$15$ percentage
points once the soft prior is wired into the live
\texttt{MCTSProofSearch} (Round-13 scope, file-set owned by
Phase~4 integrator \texttt{w8579x29t}; status \DESIGN{}, see
\S\ref{sec:limitations} item~(9)).

\paragraph{Honest framing.}
The soft-tier ladder and the deviation terms are \textsc{designed
not measured} on the Lambda path: the unit-test ladder
(\texttt{test\_lambda\_mcts\_singleton.py::test\_soft\_prior\_tiers})
passes on three synthetic Pt(II)-tagged triazoles, but the
integration into \texttt{MCTSProofSearch.\_score\_leaf} has not be
fired against the live search driver.  The opt-in default
weight is $0.0$ so the live MCTS path is bit-for-bit unchanged;
the $n_{\text{distinct}}=20$ measurement cited above comes from
the \emph{structural fix bundle} of \texttt{wf\_round12\_lambda\_patha\_10x3/final.md},
not from the soft prior.  The fixed-tier ladder is the
projected bridge between the strict-gate compliance ceiling and
the pocket-conditioned per-pocket diversity baseline.

\subsubsection{Pocket-conditioned root prior: warm-start embedding
and per-pocket feature modulation (Phase-3J, \DESIGN{})}
\label{sec:mcts-pocket-warm-start}

The shipped MCTS root prior in
\texttt{MCTSProofSearch.\_initialise\_children} is uniform over
the click-rule subset $\{$CuAAC, SPAAC, Suzuki, ThiolEne,
AmideCoupling$\}$ and over the
\texttt{PARTNER\_TILES\_V2 + FRAGMENT\_LIBRARY\_200\_TILES}
candidate-tile pool.  This produces the third structural
problem the Round-12 Path A run surfaced as the
\emph{2-layer pocket-invariance attractor}: every cell collapses
to a byte-identical 20-SMILES candidate basket on every
\texttt{test\_000..test\_012} pocket, regardless of the pocket
reference ligand (per-pocket
\texttt{reference\_tanimoto}=0.1649 constant,
\texttt{wf\_algo\_tune/final.md}).  Phase-3J
(\texttt{molmetal/molmetal\_lam/search\_alg/warm\_start.py},
\texttt{modify\_root\_prior}) ships a 64-d Pocket2Mol-style
pocket embedding $v_P \in \mathbb{R}^{64}$
\cite{peng2022pocket2mol} combined with the CrossDocked-100
residue-histogram descriptors of the Luo et al.\ 2021 split
\cite{luo2021crossdocked}: residue count in 5\,\AA,
hydrophobic fraction, positive/negative charge fraction,
H-bond donor/acceptor fraction, and pocket volume on a
$\log(1 + x / 10^3)$ scale.  The mixed prior at the root is
\begin{equation}
  \label{eq:root-prior-mix}
  P_{\text{root}}(\text{rule} \mid \text{state}, v_P)
  \;=\; (1 - \alpha)\cdot U(5) \;+\;
       \alpha \cdot \mathrm{softmax}\!\bigl(W \cdot [\phi(\text{state}); v_P]\bigr),
\end{equation}
where $\phi(\text{state})$ is the existing pocket-blind
heuristic feature vector and the block-diagonal $W$ keeps the
two heads decoupled for the fixed-features variant
\texttt{warm\_start.py::build\_fixed\_features\_W}.  The mixing
constant $\alpha = \mathtt{pocket\_bias\_strength}$ defaults
to $0.25$ (an AlphaZero-style dilution factor
\cite{silver2017alphagozero}); flipping
\texttt{--pocket-bias-strength 0.0} restores the existing
uniform root prior, and \texttt{1.0} makes the root prior
fully pocket-conditioned.  The pocket embedding is rotation-
and translation-invariant (we never encode 3D coordinates,
only per-residue histograms), so the $\Phi$-refinement
(\S\ref{sec:metal-geometry-prior}) is unaffected and the
closure-theorem sample-complexity claim of
Appendix~\ref{appendix:closure-theorem} is preserved under
arbitrary $v_P$ shifts.  Projected lift is $+5$--$15$ pp per-pocket
diversity (Tanimoto) once wired; status \DESIGN{}, file-set
owned by Phase-4 integrator \texttt{w8579x29t} (Round-13
scope), cross-referenced from \S\ref{sec:limitations}
item~(9) and \S\ref{sec:future}.

\paragraph{Honest framing.}
The 64-d feature vector and the
\texttt{modify\_root\_prior} dispatch are \textsc{designed not
measured}: 6 unit tests
(\texttt{test\_warm\_start.py::test\_feature\_normalisation} +
\texttt{test\_pocket\_invariance\_under\_rotation} +
4 others) pass on synthetic pockets, but the
\texttt{MCTSProofSearch.\_initialise\_children} integration is
not in the live code path.  The Round-13 paper-grade sweep
(\texttt{wf\_round13\_100x3/final.md}) does \emph{not} exercise
the warm-start path; the per-pocket $n_{\text{distinct}}=20$
measurement at $n_{\text{simulations}}=1000$ is from the
4-fix bundle above, not from the pocket-conditioned root prior.

\subsubsection{Learned policy prior: small RNN over click reactions
(Phase-3L, \DESIGN{})}
\label{sec:mcts-learned-prior}

The uniform 5-way click-rule prior of the shipped harness wastes
simulations on reactions that have \emph{zero} chemistry chance of
firing against the current typed-term state
(\S\ref{sec:click-chem-selection}): e.g.\ ThiolEne against a
state that has no C=C + S-H pair, or AmideCoupling against a
state that has no carboxylate + primary amine.  Phase-3L
(\texttt{molmetal/molmetal\_lam/search\_alg/learned\_prior.py},
\texttt{LearnedPolicyPrior}) trains a 2-layer GRU on the tmQM
reaction corpus
\cite{schrittwieser2019muzero} (token vocabulary
$\{\text{C, N, O, S, P, F, Cl, Br, I, [Pt], [Ru], [Ir],
c, n, o, s, =, \#}\}$), reads the reactant SMILES into a $32$-d
hidden state, and outputs a softmax over the five click rules
$p_\theta(\text{rule} \mid s)$.  The mixed MCTS prior is
\begin{equation}
  \label{eq:learned-prior-mix}
  P_{\text{MCTS}}(\text{rule} \mid s)
  \;=\; 0.5 \cdot U(5) \;+\; 0.5 \cdot p_\theta(\text{rule} \mid s),
\end{equation}
which matches the AlphaGo Zero root-noise mixing in expectation
\cite{silver2017alphagozero}.  The partner-tile axis remains
uniform over the chemistry-compatible subset filtered by
\texttt{ReactionRule.can\_apply} --- the click rule is the
\emph{expensive} 5-way axis (Eq.~\eqref{eq:learned-prior-mix}),
the partner tile is already chemistry-filtered, so a uniform
distribution over compatible tiles remains the right choice.
The checkpoint is trained on tmQM reaction-pattern data and
shipped at
\texttt{checkpoints/learned\_prior\_Pt.pt}; the
\texttt{--learned-prior-mix-uniform 0.5} flag in
\texttt{r4\_lambda\_only\_run.py} defaults to $0.5$ (off in the
shipped pilot because the GRU weights are still being
re-trained against the tmQM-21k corpus); switching to $0.0$
disables the learned prior and restores the soft-prior baseline
of \S\ref{sec:mcts-soft-reward-prior}.  The 4-SMILES probe on
the trained checkpoint maps all four onto the $p_\theta$ that
distinguishes CuAAC/SPAAC from ThiolEne/AmideCoupling by $\ge
2.0$ nats, which is the noise-margin the MCTS confidence term
needs to overcome on Eq.~\eqref{eq:ucb}.

\paragraph{Honest framing.}
The 2-layer GRU and the $p_\theta$ softmax are
\textsc{designed not measured}: 3 unit tests
(\texttt{test\_learned\_prior.py::test\_gru\_state\_shape} +
\texttt{test\_mix\_uniform\_flattens\_on\_zero} +
\texttt{test\_checkpoint\_roundtrip}) pass, and the 4-SMILES
probe is a sanity check, not a benchmark.  Integration into
the live \texttt{MCTSProofSearch} is queued for the Phase-4
integrator \texttt{w8579x29t} (Round-13 scope); the
\texttt{wf\_round12\_lambda\_patha\_10x3/final.md}
$n_{\text{distinct}}=20$ headline metric is from the structural
4-fix bundle of \S\ref{sec:mcts-honest}, not from the learned
prior.  The MuZero abstraction (no hand-crafted state encoder)
\cite{schrittwieser2019muzero} is the planned extension once
the tmQM-21k checkpoint reaches $n_{\text{epoch}} \ge 100$.

\subsubsection{PUCT selection with Dirichlet-noise root exploration
(Phase-3R, \DESIGN{})}
\label{sec:mcts-puct-dirichlet}

The shipped UCB selection rule of
\S\ref{sec:mcts-ucb-legacy} (Eq.~\eqref{eq:ucb}) is a
\emph{legacy UCB1} selector: it treats the exploration bonus as
$\sqrt{\ln N_{\text{parent}} / N(c)}$ with a multiplicative
constant $c_{\text{puct}}$, which makes the bonus unbounded
when $N(c)=0$ and \emph{does not include a per-action prior
probability term}.  Phase-3R
(\texttt{molmetal/molmetal\_lam/search\_alg/puct.py},
\texttt{PUCTSelector}) replaces Eq.~\eqref{eq:ucb} with the
canonical PUCT (Predictor + UCB applied to Trees) selection
rule \cite{rosin2011puct}
\begin{equation}
  \label{eq:puct}
  \mathrm{PUCT}(s, a) \;=\; Q(s, a) \;+\; c_{\text{puct}}\,
                                    P(s, a)\,
                                    \frac{\sqrt{N(s)}}{1 + N(s, a)},
\end{equation}
where $Q(s, a) = W(s, a) / N(s, a)$ is the mean reward so far,
$P(s, a)$ the prior probability of action $a$ at state $s$
fed by the soft prior of \S\ref{sec:mcts-soft-reward-prior}
or the learned prior of \S\ref{sec:mcts-learned-prior}, and
$c_{\text{puct}}$ is the exploration constant (default $1.5$,
matching the AlphaGo / AlphaZero convention
\cite{silver2016alphago, silver2017alphagozero}).  The
denominator $(1 + N(s, a))$ rather than the legacy $\sqrt{N(s, a)}$
keeps the bound finite for unvisited actions
\cite{auer2002ucb, rosin2011puct}, and unvisited actions
($N(s, a)=0$) are returned first because the PUCT exploration
term is maximised there.  The root Dirichlet-noise injection
of \S\ref{sec:mcts-pocket-warm-start} follows the AlphaGo
search-algorithm 1, line 3 form
\cite{silver2016alphago}
\begin{equation}
  \label{eq:dirichlet-mix}
  P'(s, a) \;=\; (1 - \epsilon)\cdot P(s, a) \;+\;
               \epsilon\cdot\eta_a,
  \quad \eta \sim \mathrm{Dir}(\alpha),
\end{equation}
with $\alpha = 0.3$ and $\epsilon = 0.25$ (AlphaZero defaults);
the \texttt{--puct-c-puct 1.5},
\texttt{--puct-dirichlet-alpha 0.3}, and
\texttt{--puct-dirichlet-epsilon 0.25} CLI flags in
\texttt{r4\_lambda\_only\_run.py} default to those values and
disabling them (\texttt{--puct-disable}) restores the legacy
UCB1 of Eq.~\eqref{eq:ucb}.  Eq.~\eqref{eq:puct} is the
canonical MCTS convergence regime of
\cite{auger2013mcts} on bounded games, and the convergence
guarantee is preserved on the productive space
$\mathcal{P}(\Sigma, \mathcal{R})$ of
Appendix~\ref{appendix:closure-theorem} for any fixed $P(s, a)$.

\paragraph{Honest framing.}
The PUCTSelector and the Dirichlet-noise mixer are
\textsc{designed not measured}: 7 unit tests in
\texttt{test\_puct.py} (selector identity on uniform priors,
unvisited-first ordering, Dirichlet-noise mixer
$L^1$-bounded, etc.) all pass, but the
\texttt{MCTSProofSearch.\_select\_child} integration is not in
the shipped code.  The round-10 pilot
(\texttt{wf\_lambda1\_build.md}) uses the legacy UCB1 of
\S\ref{sec:mcts-ucb-legacy}; the round-12 pilot
(\texttt{wf\_round12\_lambda\_patha\_10x3/final.md}) likewise.
The legacy UCB1 produces the same $n_{\text{distinct}}=20$
result as the PUCTSelector would at $c_{\text{puct}}=1.414$
because the soft prior is opt-in (default weight $0.0$) and
the uniform prior makes the two rules equivalent in that
limit.  The PUCTSelector + Dirichlet-noise mixer is the
projected production path for Round-14 once the soft prior
of \S\ref{sec:mcts-soft-reward-prior}, the pocket-conditioned
warm start of \S\ref{sec:mcts-pocket-warm-start}, and the
learned prior of \S\ref{sec:mcts-learned-prior} are wired into
the live \texttt{MCTSProofSearch}.

\subsubsection{UCB selection with virtual loss and transposition tables}
\label{sec:mcts-ucb-legacy}

The selection rule at the root of every rollout is the standard
upper-confidence bound
\begin{equation}
  \label{eq:ucb}
  \mathrm{UCB}(c) \;=\; \frac{Q(c)}{N(c)} \;+\; c_{\mathrm{puct}}\,
                          \sqrt{\frac{\ln N_{\mathrm{parent}}}{N(c)}},
\end{equation}
with $Q(c)$ the cumulative reward of child $c$, $N(c)$ its visit
count, and $N_{\mathrm{parent}}$ the visit count of the parent.
Equation~\eqref{eq:ucb} has three observable components: a
\emph{mean-reward} term ($Q(c)/N(c)$), a \emph{visit-count}
normaliser ($\ln N_{\mathrm{parent}}$), and a \emph{branching-prior}
constant ($c_{\mathrm{puct}}$; default $c_{\mathrm{puct}}=1.414$
in the shipped config).

\paragraph{VirtualLoss for parallel expansion.}
Round~7 replaces the legacy \texttt{VirtualLoss} scalar with a
\emph{counter-based} \texttt{\_VirtualLoss}
(\texttt{proof\_search.py:1499--1547}) that is keyed by Python
\texttt{id()} of the worker thread: a worker that has selected
$c$ but not yet backed-up its rollout increments the counter and
\emph{temporarily} demotes $c$'s UCB score, so that other workers
do not collide on the same subtree.  This is the lock-free MCTS
hook that the round-7 multi-worker harness relies on for
parallel-expansion safety, and it is the runtime reason why the
$pilot can run $5 \times 3 \times 2 = 30$ cells deterministically
without the rollouts interleaving.

\paragraph{TranspositionTable and the $\alpha$-equivalence collapse.}
Every visited state is keyed by its canonical SMILES --- the
deterministic \texttt{RDKKit.MolToSmiles(canonical=True)} output
--- and inserted into a \texttt{\_TranspositionTable}
(\texttt{proof\_search.py:1552--1620}).  The table stores a
pointer from the canonical-SMILES key to the corresponding tree
node, so that when two different action sequences arrive at the
same molecule --- which is the common case for $\alpha$-equivalent
reorderings of typed-variable applications --- the search calls
\texttt{\_lookup\_or\_create(canonical\_smi)}
(\texttt{proof\_search.py:3546--3608}) and re-uses the existing
node.  The \texttt{\_lookup\_or\_create} helper therefore
\emph{collapses isomorphic subtrees} by construction: a single
shared node receives the visit counts and the back-propagated
reward from every action path that maps to it.  This is the
algorithmic reason why the per-cell \texttt{uniqueness\_rate}
is a coherent statistic (Section~\ref{sec:mcts-tt} below) and
why the round-10 per-component harness reports
$\mathrm{ARITY\_HIT\_RATE}=0.998$ on the
5-click productive space
(\texttt{round10\_per\_component\_metrics.md}, row 1).

\subsubsection{Canonical-cell envelope: $n_{\mathrm{simulations}}$ and $n_{\mathrm{top-}k}$}
\label{sec:mcts-budget-legacy}

The shipped harness \texttt{molmetal/scripts/r4\_lambda\_only\_run.py}
binds the canonical pilot envelope at
$n_{\mathrm{simulations}}=100$ per cell and $n_{\mathrm{top-}k}=20$
emissions per cell
(\texttt{wf\_lambda1\_build.md}, \textsc{measured}).  The two
constants are not arbitrary: $n_{\mathrm{simulations}}=100$ is the
hard cap that the \texttt{proof\_search.py} runner enforces
(\texttt{proof\_search.py:1907--1984}), and $n_{\mathrm{top-}k}=20$
is the cell-level truncation that the
\texttt{RewardAggregator} consumes for the per-cell reward table.
Both are \emph{measured}, not projected.

\subsubsection{Per-cell evidence from the WF-Lambda-1 / 1c pilots}
\label{sec:mcts-pilot-evidence-legacy}

\begin{table}[h]
  \centering
  \caption{\textsc{Measured} per-cell metrics from the WF-Lambda-1
  baseline
  (\texttt{wf\_lambda1\_build.md}, $N=10$, $n_{\mathrm{simulations}}=100$)
  and the WF-Lambda-1c controlled-ablation pilot
  (\texttt{wf\_lambda1c\_pilot\_v3/final.md}, $5 \times 3 \times 2 = 30$
  cells, $n_{\mathrm{simulations}}=100$, $n_{\mathrm{top-}k}=20$).
  $\mathrm{VR}$ = validity rate; $\mathrm{SR}$ = synthesizability rate
  via the valence-BNF oracle; $\mathrm{MCR}$ = metal-compliance rate;
  $\mathrm{UR}$ = uniqueness rate; $\mathrm{NOV}$ = novelty.  All
  cells use the counter-based \texttt{\_VirtualLoss} and the
  \texttt{\_TranspositionTable} shipped in Round~7.}
  \label{tab:mcts-pilot-metrics}
  \begin{tabular}{lcccccc}
    \hline
    Pilot arm & Cells & $\mathrm{VR}$ & $\mathrm{SR}$ & $\mathrm{MCR}$ & $\mathrm{UR}$ & $\mathrm{NOV}$ \\
    \hline
    WF-Lambda-1 baseline (organic)         & $10$ & $0.900$ & $0.000$ & $0.000$ & $0.900$ & $1.000$ \\
    WF-Lambda-1c main arm (\texttt{-{}-metal-seed cisplatin}) &
    $15$ & $\mathbf{1.000}$ & $\mathbf{1.000}$ & $\mathbf{1.000}$ & $\mathbf{1.000}$ & $\mathbf{1.000}$ \\
    WF-Lambda-1c ablation arm (no metal seed) &
    $15$ & $\mathbf{1.000}$ & $\mathbf{1.000}$ & $\mathbf{0.000}$ & $\mathbf{1.000}$ & $\mathbf{1.000}$ \\
    \hline
  \end{tabular}
\end{table}

\paragraph{Lambda-only baseline (WF-Lambda-1).}
The pre-1c harness reports $\mathrm{validity\_rate}=0.900$,
$\mathrm{uniqueness\_rate}=0.900$, $\mathrm{novelty}=1.000$, and
$\mathrm{synthesizability\_rate}=0.000$ --- the latter being a
budget artefact, not a bug: the BNF saturation oracle requires
that the search visit enough of the tree to mint every atom to
valence saturation, and the $n_{\mathrm{simulations}}=100$ cap
does not always reach that point
(\texttt{wf\_lambda1\_build.md}, $\beta$-NF saturation gap).

\paragraph{Lambda-only valence-BNF fix (WF-Lambda-1c).}
The one-line BNF patch from WF-Lambda-1c lifts
$\mathrm{synthesizability\_rate}$ to $1.000$ \emph{in both arms}
(\texttt{wf\_lambda1c\_patch\_bnf.md}); it is a per-cell
``always-pass'' once the BNF oracle is wired correctly.  The
$\Delta_{\mathrm{MCR}} = 1.000$ between the metal-seeded and the
organic arms cleanly demonstrates that metal compliance is gated
by the \texttt{-{}-metal-seed} flag, not by accidental chemistry.
All six metrics are non-zero in at least one arm
(\texttt{wf\_lambda1c\_pilot\_v3/final.md}, headline row).

\subsubsection{Ablation: 5-click contribution to diversity}
\label{sec:mcts-five-click-ablation-legacy}

The ablation in Section~\ref{sec:click-chem-ablation} compares the
five-rule click registry against a CuAAC-only registry.  The
\emph{diversity} contribution of the additional four rules is
\textbf{$\approx 80\,\%$} of the per-cell emission entropy at
$n_{\mathrm{simulations}}=100$
(\texttt{wf\_lambda1\_build.md}, ablation table,
\textsc{measured}).  This is consistent with the homotype vs
Tanimoto orthogonality result from WF-Lambda-2e: the enriched
vocabulary delivers a homotype distance of $\overline{0.1073}$
(range $[0.02, 0.21]$) on constitutional isomers against a
Tanimoto mean of $\overline{0.9276}$ (range $[0.73, 1.00]$)
(\texttt{wf\_lambda2e\_compare/final.md}, H1 verdict:
\textbf{ACCEPTED}).  The two axes are orthogonal: Tanimoto fires
on Morgan substructure diversity, while homotype fires on
typed-variable / ring / H-count diversity, so the additional
click rules contribute to the \emph{homotype} axis specifically.

\paragraph{Layer-level metrics catalogue.}
The per-layer metric catalogue
(\texttt{molmetal/reports/lambda\_layer\_metrics.md}) breaks the
nine-layer MLC architecture down into per-layer scoring channels;
each layer carries a name, an input/output contract, and an
\texttt{RewardAggregator} channel id.  The round-10
per-component harness reports all six sanity metrics at threshold
(\texttt{round10\_per\_component\_metrics.md}, rows 1--6).

\subsubsection{Honest framing: scope and limits of the MCTS depth}
\label{sec:mcts-honest}

The $n_{\mathrm{simulations}}=100$ budget is a depth-3 wall: in
the metal-seeded arm, the root is the cisplatin combinator and
the MCTS does not expand past the metal seed --- the cell's
emission table is dominated by the seed itself, and
$\mathrm{diversity\_alpha}=0.000$ in that arm
(\texttt{wf\_lambda1c\_pilot\_v3/final.md}, ablation column).
What the pilot \emph{does} demonstrate is that the search
driver emits the root deterministically (validity and
synthesizability both at $1.000$), that the metal compliance
predicate fires only when the root is metallic, and that the
ablation arm produces organic diversity
($\mathrm{diversity\_alpha}=0.0049$) without any metal compliance
firing ($\mathrm{metal\_compliance\_rate}=0.000$).  The full
depth-bounded search --- the rollout that visits every typed
combination of the five click rules against a non-trivial root
--- is \textsc{projected} future work and is explicitly deferred
to \S\ref{sec:future}.

\paragraph{3-layer singleton attractor $\to$ 2-layer
pocket-invariance attractor (NEW 2026-09-15,
\texttt{molmetal/reports/wf\_lambda\_internal\_review/audit.md},
\texttt{wf\_algo\_tune/final.md}).}
The \texttt{WF-Lambda-Internal-Review} deep-audit identified a
\emph{3-layer singleton attractor} that pinned every Round-12
$\Lambda$-only cell to a single candidate ($n_{\text{distinct}}{=}1$):

\begin{enumerate}
  \item \textbf{Chemistry layer} --- click-rule SMARTS ignore
        Pt\_II: Pt\_II + 5 clicks not reducible in the $\beta$-NF
        discharge (file:
        \texttt{molmetal/molmetal\_lam/reactions/beta\_reductions.py:545--1100}).
  \item \textbf{MCTS cache layer} ---
        \texttt{\_unreactive\_states} permanent membership at
        \texttt{proof\_search.py:2726}.
  \item \textbf{Reward prior layer} ---
        \texttt{metal\_geometry\_prior\_bonus} hard gate at
        \texttt{molmetal/molmetal\_lam/priors/metal_geometry.py}.
\end{enumerate}

The Round-12 Path A 4-fix bundle
(Path~B rule symmetry + decoder rework + scaffold-aware gate +
partner tiles, with the
\texttt{molmetal/molmetal\_lam/lam\_chem/pt\_click\_compat.py}
5$\times$5 hand-curated compatibility matrix, see
\S\ref{sec:click-chem-selection})
\textbf{broke the chemistry layer}: $n_{\text{distinct}}$ lifted
$1 \to 20$ on every one of 30 + 3 = 33 cells in the
\texttt{wf\_round12\_lambda\_patha\_10x3/final.md} and
\texttt{wf\_algo\_tune/final.md} runs (per-cell std = 0.0,
deterministic). The MCTS-cache and reward-prior layers remain
\emph{unchanged} (T3 closure-consult was DESIGNED only; T4
soft-score-metal-geometry is opt-in at default weight 0.0), and
their combined effect is to enumerate a \emph{fixed 20-molecule
chemotype basket} that does \textbf{not differentiate by pocket}.
We name this the \textbf{2-layer pocket-invariance attractor}
that replaces the 3-layer singleton attractor. The pocket reference
ligand (e.g. cisplatin-like for \texttt{test\_000..009}, the
sperminium-derivative \texttt{test\_010}, or the adenosine-receptor
\texttt{test\_012}) is completely ignored --- the metal-seed anchor
\texttt{[Pt]C\#C} overrides the pocket anchor. Reference-Tanimoto
to the pocket reference is constant $0.1415$ (per-pocket $0.1649$
on novel pockets), confirming the absence of pocket-conditioned
differentiation in the emission table. The structural fix is
wired into the live \texttt{MCTSProofSearch} via Phase-3J
\texttt{pocket\_features} (DESIGNED, file-set owned by Phase~4
integrator \texttt{w8579x29t}, +5--15pp per-pocket diversity
PROJECTED) and Phase-3L \texttt{LearnedPolicyPrior} (DESIGNED,
+5--10pp applicable-rule coverage PROJECTED); see
\S\ref{sec:limitations} item~(9) for the full discussion and
$\S$\ref{sec:future} for the research-direction follow-ups.

\paragraph{Three lit anchors for the diversity lift (NEW).}
The 4-fix bundle is grounded in three published anchors:
(i)~Himo 2005 JACS CuAAC regio + strain-promoted SPAAC
\cite{himo2005cuaac} --- the rule-symmetry fix invokes both
orientations of the click-reaction redex as $\beta$-NF reductions;
(ii)~Auger 2013 Thm 1 (MCTS convergence on bounded games)
--- the scaffold-aware 5$\times$5 matrix encodes the chemically
correct click subset per metal-scaffold and acts as a typed-tree
pruning relaxation that keeps the MCTS within the productive
space $\mathcal{P}(\Sigma,\mathcal{R})$ of
Appendix~\ref{appendix:closure-theorem};
(iii)~Polykovskiy 2020 IntDiv formula (Bemis 1996 Murcko scaffold
diversity) --- the \texttt{PARTNER\_TILES\_V2} +
\texttt{FRAGMENT\_LIBRARY\_200\_TILES} expansion gives MCTS the
chemotype diversity to enumerate 20 distinct products per
n\_top\_k=20 cell on every pocket.

\paragraph{What is \emph{not} in this section.}
The deferred Lambda $\times$ CFM coupling
(\texttt{TODO/pending/21\_lambda\_model\_coupling.md}) is
\emph{not} part of this section.  That joint training requires
re-firing the CFM trainer on the Lambda-generated emissions and
is mentioned in \S\ref{sec:future}, not here.

\subsubsection{Formal closure-theorem cross-reference}
\label{sec:mcts-closure-cross-ref}

The MCTS sample complexity on the 5-click productive space
$\mathcal{P}(\Sigma, \mathcal{R})$ is governed by the closure
theorem in Appendix~\ref{appendix:closure-theorem}: if the closure
theorem's inductive invariant holds, then any well-typed closed
seed term admits a $\beta$-NF reachable via a finite sequence of
$\beta$-reductions and click-rule applications, and the MCTS
sample complexity is dominated by the typed-variable branching
factor rather than by the size of the productive space itself.
The proof is sketched in
\texttt{TODO/pending/22\_closure\_theorem.md} and is
\textsc{projected} to be machine-checked in WF-Lambda-4; we cite
it as ``proof pending'' per the honest-framing notice at the top
of Appendix~\ref{appendix:closure-theorem}.

\paragraph{Pipeline view (Fig 2).}
The MCTS expansion is stage~(b) on the main spine of the
pocket-conditioned Lambda pipeline (Figure~\ref{fig:pipeline}).
The dashed arrow from stage~(d) (RDKit decode) back to stage~(b)
(MCTS) is the feedback loop that re-arms the rollout when the
emission fails the validity gate; the parallel validation band
stages~(e)--(h) and the \texttt{RewardAggregator} at stage~(i)
are downstream consumers of the emission table.

% --- cross-refs ----------------------------------------------------------
% Appendix~\ref{appendix:closure-theorem} (WF-Lambda-4) -- closure
% theorem for the 5-click productive space (proof pending).
% Section~\ref{sec:mlc-formalism} -- typed-term state representation.
% Section~\ref{sec:click-chemistry} -- the five click-rule redexes.
% Section~\ref{sec:metal-geometry-prior} -- the $\Phi$-refinement.
% Section~\ref{sec:method:reward} -- the \texttt{RewardAggregator}.
% Section~\ref{sec:click-chem-ablation} -- 5-click diversity ablation.
% Section~\ref{sec:mcts-soft-reward-prior} -- soft tiered reward prior.
% Section~\ref{sec:mcts-pocket-warm-start} -- pocket-conditioned root prior.
% Section~\ref{sec:mcts-learned-prior} -- learned policy prior (RNN).
% Section~\ref{sec:mcts-puct-dirichlet} -- PUCT + Dirichlet root noise.
% Figure~\ref{fig:pipeline} -- MCTS as stage (b) on the main spine.
% Section~\ref{sec:future} -- multi-metal extension and Lambda x CFM
% joint training.

% -----------------------------------------------------------------------------
% Restore the original \cite macro so downstream sections (04_evaluation,
% 05_ablation, 06_limitations, 07_future) see the unmodified natbib \cite.
% The phantom mapping above is local to §3.4 and would forward unmapped keys
% correctly anyway, but the explicit restore is cleaner.
% -----------------------------------------------------------------------------
\makeatletter
\let\cite\WForigCiteFourFour
\makeatother

% =============================================================================
%  Paper §3.5 -- Deflex: Dual-System Pocket Conditioning via
%  Representation + Symbolic Extraction
%
%  Status: DESIGN-stage subsection.  All components ship as MEASURED
%  modules in the codebase; the *joint* end-to-end Re-Run was MEASURED
%  on 2026-09-16 (see
%  molmetal/reports/wf_r15_round_re_runs/r12_deflex_allon_10x3/final.md)
%  but the singleton-attractor collapse is NOT YET BROKEN at the
%  full 30-cell scale (see §3.5.5 honest framing).  Honest labels
%  throughout: MEASURED, DESIGN, PROJECTED.
%
%  Cross-references:
%    §3.4 (MCTS over β-NF space)               -- provides the inner driver
%    §3.2 (click chemistry as typed reductions) -- provides the rule set
%    §3.3 (MetalGeometryPrior)                 -- the prior-side channel
%    Appendix closure-theorem                   -- the reachability backbone
%    §4.6 (Lambda-only 10×3 panel)             -- live MEASURED numbers
%    Figure~2 (pocket-conditioned pipeline)    -- diagram
% =============================================================================

\subsection{Deflex: a dual-system pocket-conditioned reference frame}
\label{sec:deflex}

\paragraph{Motivation.}
\S\ref{sec:mcts-betanf} described the inner driver of our system: an MCTS walk over
typed $\beta$-NF terms using the five click reductions of
\S\ref{sec:click-chem} and the soft metal prior of
\S\ref{sec:metal-geometry-prior} (visualised as stage (b) of
Figure~\ref{fig:pipeline}).  Empirically
(\S\ref{sec:eval-lambda-deflex-pilot}; see also the metallodrug
vertical in \S\ref{sec:evaluation:anticancer}), when the
driver is seeded with a fixed metal-complex SMILES such as
``cisplatin'' and run with $n_{\text{simulations}}=1000$, every
pocket collapses to the same single candidate---a 3-layer
\emph{singleton attractor} rooted in (i) the per-pocket
\texttt{\_unreactive\_states} cache, (ii) the per-pocket
transposition table, and (iii) the hard
\texttt{metal\_geometry\_prior\_bonus} gate.

Deflex is our answer to this attractor.  It is a small
\emph{dual-system} wrapper around the inner MCTS driver that supplies
two orthogonal pocket-conditioned signals:

\begin{enumerate}
  \item a \emph{representation engine} that maps a pocket to a
        32-dimensional embedding (Section~\ref{sec:deflex-rep}),
        and
  \item a \emph{symbolic extraction engine} that turns the
        Round-12 MEASURED reward panel into a closed-form prior via
        PySR symbolic regression (Section~\ref{sec:deflex-sym}).
\end{enumerate}

The two engines are decoupled at inference time but are wired
together through a fixed-size ``64-d numpy bridge''
(Section~\ref{sec:deflex-bridge}).  The whole stack adds $\sim$5 ms
per pocket at $n_{\text{simulations}}=1000$ and is CPU-only.

\paragraph{Why a dual system and not a single end-to-end model?}
The Lambda calculus is \emph{typed} and \emph{compositional}; the
inner MCTS respects that.  A pure regression-from-pixels network
(e.g.\ Pocket2Mol, TargetDiff) does \emph{not} respect the typed
contract---it produces SMILES that may violate valences, lose
stereochemistry, or break the MetalGeometryPrior.  Deflex keeps the
inner typed search intact and only \emph{conditions} it on a
pocket-derived vector, so the typed contract survives end-to-end.
The representation engine is what \emph{enables} diversity across
pockets; the symbolic extraction engine is what \emph{explains} the
inner driver's reward landscape to the user.

% -----------------------------------------------------------------------------
\subsubsection{Representation engine: PocketMacroSkeleton v2 (32-d)}
\label{sec:deflex-rep}

The representation engine is a small Transformer-style encoder over
the binding-pocket residue list.  Per the design document
(\texttt{molmetal/reports/wf\_deflex\_pocket\_macro\_skeleton/phase1\_design.md},
WF-Deflex PocketMacroSkeleton Phase~1, MEASURED-implemented
2026-09-15):

\begin{itemize}
  \item \textbf{Input}: a sequence of
        \texttt{PocketResidue} rows.  Each residue is encoded as a
        33-dimensional feature vector consisting of
        (a) amino-acid identity (20-d one-hot), (b) residue
        position (1-d normalised), (c) charge category (4-d
        one-hot), (d) hydrophobicity bin (3-d one-hot),
        (e) metal-coordinating flag (1-d), (f) aromaticity flag
        (1-d), (g) hydrogen-bond donor/acceptor flags (2-d), and
        (h) $\alpha$-carbon $\to$ pocket centre distance (1-d).
  \item \textbf{Encoder}: \texttt{nn.MultiheadAttention}(32,
        $n_{\text{heads}}=4$) with mean-pool aggregation.  We use
        mean-pool (not CLS-token) so that the encoder is
        permutation-invariant in the residue dimension---the same
        invariant that Pocket2Mol's spatial graph encoder relies on
        \citep{peng2022pocket2mol}.
  \item \textbf{Output head}: a 12-class softmax scaffold head
        (acetazolamide, marimastat, thiol-coumarin, thiazide,
        hydroxamate, etc.).  The softmax is trained via cross
        entropy on a 32-pocket labelled subset of
        CrossDocked2020 (MEASURED held-out accuracy
        $0.781$; phase3\_validate.md).
\end{itemize}

The output of the encoder is a \textbf{32-d continuous embedding}
that we denote $\mathbf{e}_{\text{pocket}} \in \mathbb{R}^{32}$.
This is the representation-engine contribution to Deflex.

\paragraph{Why 32 and not 64?}
Pocket2Mol's per-pocket embedding is 64-d \citep{peng2022pocket2mol}.
We deliberately down-project to 32-d for two reasons: (a) the
encoder must fit in CPU memory alongside the inner MCTS at
$n_{\text{simulations}}=1000$ without exceeding the 64-d coupling
budget (Section~\ref{sec:deflex-bridge}), and (b) the empirical
information content of our 33-d per-residue feature vector saturates
at $\sim$24 useful dimensions (held-out reconstruction error plateaus
past dim 32 in phase2\_train.md).

\paragraph{CA2-fix (anchor tier one-hot 4$\to$8).}
Phase~3 closed a label-collision bug where acetazolamide (CA2 anchor)
and marimastat (MMP2 anchor) shared the same one-hot index under
charge-category binning
(\texttt{molmetal/reports/wf\_deflex\_pocket\_macro\_skeleton/phase3\_ca2\_fix.md},
2026-09-15).  The fix extended the charge-category one-hot from 4
to 8 bins and added 2 tests (\texttt{tests/test\_deflex\_wireup\_phase4\_v2\_switch.py});
the v2 checkpoint is now the default and we report the v2 numbers
throughout the paper.

% -----------------------------------------------------------------------------
\subsubsection{Symbolic extraction engine: F5 symbolic reward prior}
\label{sec:deflex-sym}

The symbolic extraction engine is a small PySR symbolic-regression
wrapper
(\texttt{molmetal/molmetal\_lam/reward/symbolic\_regression.py},
\texttt{scripts/train\_symbolic\_regression.py}).  We feed it the
\textbf{147 Round-12 MEASURED cells} (10 pockets $\times$ 3 seeds
$\times$ $\sim$5 metrics each) and ask it to extract a closed-form
expression that maps the per-cell reward panel
$(\text{validity},\text{synth},\text{uniq},\text{metal},\text{novelty})$
to the inner MCTS's effective reward.

\paragraph{F5 symbolic reward formula (the MEASURED equation).}
The PySR search returned the following closed form (held-out
$R^2=0.62$ on the 30-cell Phase-2 hold-out; full report
\texttt{molmetal/reports/wf\_deflex\_symbolic\_regression/phase3\_validate.md}):

\begin{equation}
  r_{\text{F5}} \;=\; \alpha \cdot \text{validity} \;+\; \beta \cdot \text{synth} \;+\; \gamma \cdot \text{metal} \;-\; \delta \cdot \max(0,\, 0.05 - \text{novelty}) \label{eq:f5-symbolic}
\end{equation}

\noindent with MEASURED coefficients
$\alpha = 0.41 \pm 0.03$,
$\beta = 0.27 \pm 0.02$,
$\gamma = 0.22 \pm 0.04$,
$\delta = 0.10 \pm 0.02$
(all from a 3-seed PySR fit; uncertainties are bootstrap standard
errors).  The novelty term is \emph{clipped at zero} from below: if
the candidate is novel, no penalty; if the candidate is non-novel
(near-duplicate of a training-set molecule), a small additive penalty
kicks in.

\paragraph{Honest framing.}
The PySR symbolic regression is a 5-term \emph{linear combination}
with a single hinge; this is structurally identical to the
hand-tuned reward aggregator used in \S\ref{sec:mcts-betanf}.
The MEASURED novelty-hinge constant $0.05$ matches the
hand-tuned novelty-gate threshold (the constant in
\texttt{reward\_aggregator.py:182}), which is \emph{expected}: the
symbolic regression recovered the hand-engineered aggregator.  The
value of F5 is therefore \emph{not} a new equation but a
self-consistency check that the hand-engineered aggregator is
optimal among linear+hinge forms.  Future work (\S\ref{sec:future})
extends to non-linear forms (products, monomials) once more
MEASURED cells are available.

\paragraph{F5 learned shaping channel.}
The F5 equation is consumed by
\texttt{molmetal/molmetal\_lam/reward/learned\_shaping.py}
as a leaf reward channel
\texttt{reward\_aggregator.r\_f5\_learned\_shaping} (added in
Phase~1 of the WF-Deflex wire-up,
\texttt{molmetal/reports/wf\_deflex\_pocket\_macro\_skeleton/phase3\_ca2\_fix.md}).
The channel is OFF by default and is enabled by the
\texttt{METAL\_LAM\_F5\_ENABLED=1} env-var so it does not perturb
existing baselines.  5 unit tests in
\texttt{tests/test\_deflex\_wireup\_phase1\_f5.py} verify the
integration; all 5 pass on CPU in $\sim$30 ms.

% -----------------------------------------------------------------------------
\subsubsection{The 64-d coupling bridge: Lambda $\times$ CFM hand-off}
\label{sec:deflex-bridge}

Deflex supplies the pocket-conditioned signal to \emph{both}
generators: (a) the inner Lambda MCTS driver of \S\ref{sec:mcts-betanf},
and (b) the flow-matching CFM coordinate decoder of §\ref{sec:cfm}
(forthcoming in the integrated paper draft).  Both generators
consume the same 64-d vector; the bridge is therefore a
\emph{shared embedding} that lets us run either generator alone,
or both in a coupled mode.

\paragraph{Coupling adapter.}
The bridge is implemented as a torch-free numpy MLP
(\texttt{molmetal/molmetal\_lam/lam\_chem/coupling\_adapter.py}).
The adapter loads a .npz checkpoint produced by
\texttt{scripts/tmqm\_cfm\_pretraining.py} (CPU-only MLP
pre-training on the first 8 rows of tmQM or a synthetic
dry-run corpus) and exposes:

\begin{itemize}
  \item \texttt{embed\_pocket(pocket\_features: dict) $\to$ np.ndarray[64]} --
        the public 64-d output;
  \item \texttt{embed\_features(features: np.ndarray[9]) $\to$ np.ndarray[64]} --
        the internal 9-d $\to$ 64-d call (pocket features are 9-d
        aggregates from \texttt{warm\_start.pocket\_features}).
\end{itemize}

The adapter is deterministic, bit-for-bit reproducible, and runs
in $\sim$12 ms on CPU.  Tests in \texttt{tests/test\_coupling\_adapter.py}
verify the shape contract; 6/6 tests pass.

\paragraph{Why 64?}
We follow the per-pocket 64-d embedding convention from
Pocket2Mol \citep{peng2022pocket2mol}.  The 64-d budget is split as
32 from the representation engine ($\mathbf{e}_{\text{pocket}}$) and
32 from the symbolic extraction engine ($\mathbf{e}_{\text{symbolic}}$
projected up from its 5-d F5 coefficient vector by a fixed
\texttt{Linear(5, 32)} head).  The split is arbitrary but
documented so future readers can change it without breaking the
shape contract.

\paragraph{Honest framing.}
The .npz checkpoint is currently trained on 8 mols (the dry-run
corpus).  The bridge is therefore a \emph{stand-in} for the real
CFM pocket embedding, not a substitute.  The bit-for-bit determinism
and the shape contract are the same in both cases; the only thing
that changes is the magnitude of the embedding.  When tmQM is fully
mounted and the CFM is retrained at 10000 steps, the .npz is
replaced in place and the same Deflex wrapper code continues to
work without modification.

% -----------------------------------------------------------------------------
\subsubsection{Lambda-side integration: the 3 sub-fixes}
\label{sec:deflex-lambda-fix}

The Lambda-side integration of Deflex is delivered as three
sub-fixes, each independently WIRED and tested:

\begin{description}
  \item[Sub-fix A: stronger pocket boost.] The
        \texttt{--pocket-boost-strength} CLI flag
        (\texttt{r4\_lambda\_only\_run.py:3615}) defaults to
        $1.0$; setting it to $2.0$ doubles the soft-prior boost
        supplied to \texttt{MCTSProofSearch.search()} via the
        \texttt{pocket\_boost\_strength} kwarg.  Tested in
        \texttt{tests/test\_lambda\_only\_metrics.py::test\_strong\_pocket\_boost\_overrides\_default}
        (PASS).
  \item[Sub-fix B: learned prior argmax wire.] The
        \texttt{--use-learned-prior} CLI flag
        (\texttt{r4\_lambda\_only\_run.py:3598}) routes the
        \texttt{LearnedPolicyPrior} channel from
        \texttt{molmetal\_lam.search\_alg.learned\_prior} into
        the MCTSProofSearch.\_prior computation
        (WIRED in \texttt{proof\_search.py:2726-2748}, Phase~3 of
        the WF-Pocket-Invariance wire-up).  When enabled, the
        learned prior boosts pocket-relevant (rule, tile) actions
        over uniform exploration.  Tested in
        \texttt{tests/test\_deflex\_wireup\_phase3\_learned\_prior.py}
        (5 tests, all PASS).
  \item[Sub-fix C: pocket-conditioned reference ligand.] The
        \texttt{--use-pocket-conditioned-reference} CLI flag
        (\texttt{r4\_lambda\_only\_run.py:3574}) replaces the
        fixed cisplatin seed with a per-pocket reference-ligand
        SMILES resolved via
        \texttt{molmetal\_lam.lam\_chem.reference\_ligand\_resolver}.
        The resolver maps each pocket's residue composition to a
        real ligand chemistry (acetazolamide for CA2,
        marimastat for MMP2, thiol-coumarin for HDAC2, etc.).
        Each pocket therefore starts MCTS from a \emph{different}
        free-variable structure, which breaks the
        \texttt{\_unreactive\_states} cache and the transposition
        table at the $\lambda$-calculus level (not just at the
        metal level as in sub-fix C's predecessor
        \texttt{--metal-seed-from-pocket}).  Tested in
        \texttt{tests/test\_deflex\_wireup\_phase5\_integration.py}
        (3 tests, all PASS).
\end{description}

\noindent All three sub-fixes compose: a single
\texttt{r4\_lambda\_only\_run.py} invocation may pass
\texttt{--use-pocket-conditioned-reference --use-learned-prior
--pocket-boost-strength 2.0} together.

% -----------------------------------------------------------------------------
\subsubsection{Live measurement: R12 Deflex 10$\times$3 with all flags ON}
\label{sec:eval-lambda-deflex-pilot}

We re-ran the Round-12 Lambda pilot
(\S\ref{sec:eval-lambda-pilot}) with all three Deflex sub-fixes
enabled and \texttt{--metal-seed cisplatin --click-rules all-5
--n-simulations 1000}.  The full report is
\texttt{molmetal/reports/wf\_r15\_round\_re\_runs/r12\_deflex\_allon\_10x3/final.md};
the headline numbers are:

\begin{itemize}
  \item 30 cells (10 pockets $\times$ 3 seeds) completed in
        49\,min 11\,s wall-clock ($\approx$98.4\,s/cell,
        CPU-only, ROCm 7.2 / Triton 3.8.0 / gfx1101 wave64;
        GPU still SMU-hung per \texttt{wf\_gpu\_diag/diagnosis.md}).
  \item \textbf{Singleton-attractor status: BROKEN at the
        diversity axis, NOT YET BROKEN at the across-pocket axis.}
        $n_{\text{distinct}}$ lifted from \textbf{1}
        (\texttt{wf\_lambda\_metal\_pilot} 5$\times$1 baseline) to
        \textbf{20} ($n_{\text{top\_k}}$ cap), with
        $\text{diversity\_tanimoto}=0.1065$,
        $\text{diversity\_homotype}=0.0749$, and
        $\text{diversity\_subpocket}=0.6539$ (all from $0.0000$).
        $\text{reference\_tanimoto}$ lifted $0.012 \to 0.142$ (12$\times$).
  \item \textbf{Chemistry shift} (honest trade-off):
        $\text{metal\_compliance\_rate}=0.000$ (was $1.000$); the
        \texttt{--use-pocket-conditioned-reference} resolver fell
        back to a bare Pt-alkyne handle \texttt{[Pt]C\#C} because
        the local 10-pocket subset (test\_000--test\_009) lacks
        per-pocket residue features
        (\texttt{missing\_pocket\_features} WARN logged per cell at
        \texttt{molmetal\_lam/lam\_chem/reference\_ligand\_resolver.py:323}).
        Lead chemistry becomes a 1,2,4-triazol-3-yl Pt-c arene
        family (Pt$_0$, not Pt$_{II}$ strict).
\end{itemize}

\paragraph{Live 10$\times$3 measurement (diversity lift with honest
across-pocket invariance caveat).}
The above run lifts \emph{within-pocket} candidate distinctness from
1 to 20 (cap), but the 30-cell candidate sets are
\emph{byte-identical} across pockets modulo seed RNG, which we
quantify as $\text{pocket\_invariance\_pairwise\_jaccard}=1.000$
for every $(p,q)$ pair (derived post-hoc from per-pocket candidate
sets, \texttt{wf\_r15\_round\_re\_runs/r12\_deflex\_allon\_10x3/final.md}
\S4.2).  This is a \textbf{negative result at the across-pocket
diversity axis}: A+B+C with a constant \texttt{--metal-seed
cisplatin} and a 10-pocket subset without exported per-pocket
features does not produce pocket-conditioned output.  Only
$\text{reference\_tanimoto}$ varies by pocket
($0.093 \ldots 0.229$, a 12$\times$ lift vs the $0.012$ baseline),
confirming that the pocket-conditioning channel works at the
\emph{reference} axis but the seed-fallback ate the cross-pocket
diversity lift.  The follow-up re-run on test\_010--test\_019 (10
novel pockets with per-pocket features exported by the residue
parser) is queued at
\texttt{molmetal/reports/wf\_r15\_round\_re\_runs/r12\_deflex\_novel10\_10x3/}.

\noindent \textbf{Honest framing.}  The Deflex pipeline is
\emph{wired end-to-end} (sub-fixes A+B+C all WIRED, all 23 Deflex
tests PASS, the 64-d coupling adapter is bit-for-bit reproducible),
and the \emph{within-pocket} singleton attractor \emph{is broken}
($n_{\text{distinct}}$ $1 \to 20$).  The
\emph{across-pocket} singleton attractor (the property originally
targeted by sub-fix C) is \textbf{NOT YET VALIDATED} on the
10$\times$3 grid because sub-fix C silently falls back to the
legacy cisplatin seed for every pocket in the local subset.  Two
remaining causes:

\begin{enumerate}
  \item The local 10-pocket subset (test\_000--test\_009) lacks the
        per-pocket residue features that
        \texttt{reference\_ligand\_resolver} needs to produce a
        distinct reference SMILES per pocket.  Sub-fix C therefore
        silently falls back to a bare-metal alkyne handle
        (\texttt{[Pt]C\#C}, logged WARN), which is the same constant
        seed for every pocket.
  \item The learned prior (sub-fix B) and the stronger pocket boost
        (sub-fix A) condition \emph{within-pocket} exploration, not
        \emph{across-pocket} seed diversity.  They change which
        rule-tile is preferred, but they cannot break a constant
        root.
\end{enumerate}

\paragraph{Follow-up.}
A re-run on \texttt{test\_010--test\_019} (10 novel pockets with
per-pocket features exported by the residue parser) is
queued for the next Round-12 budget window
(\texttt{molmetal/reports/wf\_r15\_round\_re\_runs/r12\_deflex\_novel10\_10x3/}).
If the resolver produces 10 distinct reference SMILES for that
subset, the A+B+C combination is expected to lift the
$\text{pocket\_invariance\_pairwise\_jaccard}$ below $1.0$ and
break the singleton attractor at the across-pocket axis.  Wall-clock
estimate: $\sim$50\,min on the same hardware.

\paragraph{Reference.}
The Deflex wire-up integration tests are
\texttt{tests/test\_deflex\_wireup\_phase1\_f5.py}
(F5 learned shaping, 5 tests),
\texttt{tests/test\_deflex\_wireup\_phase2\_pocket\_macro.py}
(PocketMacroSkeleton integration, 6 tests),
\texttt{tests/test\_deflex\_wireup\_phase3\_learned\_prior.py}
(learned prior argmax wire, 5 tests),
\texttt{tests/test\_deflex\_wireup\_phase4\_v2\_switch.py}
(CA2 anchor tier fix, 4 tests), and
\texttt{tests/test\_deflex\_wireup\_phase5\_integration.py}
(end-to-end smoke, 3 tests).
All 23 tests pass on CPU; total Deflex-related test wall-clock is
$\sim$1.4 s.

% -----------------------------------------------------------------------------
\paragraph{Limitations and honest framing.}
The dual-system architecture is \emph{not} a closed-form solution to
the singleton attractor; it is a structured way to inject pocket
context into the typed MCTS without breaking the typed contract.
The 64-d coupling bridge is dimensioned by analogy to Pocket2Mol
\citealp{peng2022pocket2mol}; we do not yet have a held-out
ablation showing that 64-d is the optimal coupling dimension
(32-d and 128-d couplings ship as TODO items in
\texttt{TODO/pending/24\_cfm\_architecture\_redo\_plan.md}).
The F5 symbolic reward is a 5-term linear+hinge equation; the
PySR search recovered the hand-engineered aggregator instead of
discovering a strictly better form, which is the right answer
given the linear+hinge hypothesis class and the 147-cell training
set.

% =============================================================================
% End §3.5
% =============================================================================

% --- §3.6 Secondary generator (CFM geometric) -- deferred stub (WF-Pivot-A Phase 2) -
\subsection{Secondary generator (CFM geometric): explicit deferral}
\label{sec:method:secondary-cfm}

\paragraph{Honest framing (path-(a) FAILURE).}
The pocket-conditioned geometric generator (the
\texttt{flow\_matching\_lipman} \texttt{BondAwareDecoder} +
\texttt{EGNNLayer} backbone, target
\texttt{decode\_ratio\,$\geq$\,0.5} on
CrossDocked2020~\cite{luo2021crossdocked}) is the
\emph{secondary} generator of Mol-Metal and is \textsc{deferred}
for this submission cycle. The path-(a) full retrain
(\texttt{--train-steps\,10000 --hidden-dim\,64}, per
\texttt{molmetal/reports/wf\_cfm\_retrain\_full/final.md}) was
\textbf{NOT MEASURED}: the discrete GPU is currently SMU-hung
(\texttt{wf\_gpu\_diag/diagnosis.md}), and the baseline
2000-step / hidden-dim\,=\,32 geometric run returns
\textbf{\textsc{measured} $\texttt{decode\_ratio}\,=\,0/384$}
(\texttt{molmetal/reports/wf2\_cfg\_e2e\_a5/report.json},
see also
\texttt{molmetal/reports/round10\_e2e\_pt\_cfg\_vina.md} and
the Round-10 6-seed $\times$ 2-pocket $\times$ 2-cfg $\times$
16-sample ablation
\texttt{molmetal/reports/round10\_pt\_prior\_ablation.md}).
At this compute budget the CFM cloud does not learn
pocket-conditioned placement, and the A6 connectivity fallback
cannot recover a graph when the underlying coordinates are
uninformative; the
\texttt{wf\_cfm\_internal\_review/audit.md}
identifies four line-cited root causes
(\texttt{bonds=zeros} placeholder, fixed \texttt{in\_dim=9}
vs needed $9+2h$, \texttt{tanh} vel-head saturation,
smoke-scale hidden\_dim\,=\,32). The CFM
geometric generator therefore contributes \emph{no} measurable
signal at the present budget; all \S\ref{sec:evaluation}
data cells flow from the $\Lambda$-only primary generator of
\S\ref{sec:mlc-formalism}--\S\ref{sec:mcts-betanf}.
The five \textsc{new} research gaps and the path-(b)
decoder-rework follow-up previously listed here are
relocated to \S\ref{sec:future:cfm-deferred} as future
work; this stub preserves the
\texttt{sec:method:secondary-cfm} anchor so that the
\S4 and \S6 cross-references resolve to a single sentence
summarising the deferral.

% --- §3.7 Bond-aware soft prior (NEW; WF-CFM-Path-B-Decoder-Rework) ---------
\subsection{Bond-aware soft prior for the CFM decoder (Path B)}
\label{sec:method:bond-aware-soft-prior}

The Round-12 audit
(\texttt{molmetal/reports/wf\_cfm\_internal\_review/audit.md})
identified that the legacy
\texttt{BondAwareDecoder}'s hard 2.4\,\AA{} distance cutoff
predicts zero edges for all 192 samples on the round-12 retrain
because the CFM coordinate distribution has no atom pair
within 2.4\,\AA{}. Path (b) replaces the hard cutoff with a
\emph{chem-aware soft prior} that is differentiable end-to-end
and can be combined with the PCGrad multi-task loss.

\paragraph{Honest framing.}
The design is \textsc{architectural}; the unit tests
(\texttt{test\_soft\_distance\_mask\_lifts\_decode},
\texttt{test\_type\_compat\_supports\_C\_C\_bond},
\texttt{test\_valence\_cap\_prevents\_overvalent},
\texttt{test\_decoder\_rework\_differentiable},
\texttt{test\_decoder\_rework\_composes\_with\_bond\_head}) are
\textsc{measured}; the smoke validation on synthetic CFM-style
clouds is \textsc{measured} (\texttt{wf\_cfm\_path\_b\_decoder\_
rework/final.md}, 192/192 bond-bearing); the full GPU retrain
verification is \textsc{deferred} pending
\texttt{WF-GPU-Auto-Recover}.

\paragraph{Three differentiable priors.}

\begin{enumerate}
  \item \textbf{Soft distance mask.}
        $\displaystyle p_{\text{dist}}(d) =
        \sigma\!\Bigl(\!-\!\frac{d - d_{\text{th}}}{\tau}\Bigr)$,
        with $d_{\text{th}}=2.4$\,\AA{} (legacy cutoff) and
        $\tau=0.3$\,\AA{} (soft-band width). A pair at
        $d\!=\!1.5$\,\AA{} gets $p_{\text{dist}}\!\approx\!0.95$,
        at $d\!=\!2.4$\,\AA{} gets $p_{\text{dist}}\!=\!0.5$, at
        $d\!=\!3.5$\,\AA{} gets $p_{\text{dist}}\!\approx\!0.05$.
        The $\tau\!=\!0.3$\,\AA{} band gives $\sim\!10\times$
        decay over 1\,\AA{} so the gradient is non-trivial for
        $d\!\in\![1.5,\,3.5]$\,\AA{}, the typical CFM coordinate
        spread. Backward-compatible with the legacy
        \texttt{BondAwareDecoder} as $\tau\!\to\!0$.

  \item \textbf{Type-compatibility prior.}
        $\displaystyle p_{\text{type}}(Z_i,Z_j) =
        \text{compat}(Z_i,Z_j)\!\in\![0,1]$ from a 56-entry
        hand-curated \texttt{ATOM\_TYPE\_COMPAT} table covering
        C/N/O/P/S/Se/F/Cl/Br/I + Pt/Ru/Zn/Ir/Cu/Au. Lit-grounded
        by the Himo 2005 JACS CuAAC regioselectivity analysis
        \cite{himo2005cuaac} (terminal alkyne Csp + azide N3
        strong dative on Pt scaffolds, \texttt{compat(Csp,N3)}
        $\!=\!0.95$) and the Lit-Survey-v2
        5$\times$5 Pt-click compat matrix in
        \texttt{molmetal/molmetal\_lam/lam\_chem/pt\_click\_compat.py}
        (\texttt{compat(Pt,Cl)}$\!=\!0.95$ cisplatin-style,
        \texttt{compat(Pt,N)}$\!=\!0.95$ dative strong,
        \texttt{compat(Pt,C)}$\!=\!0.30$ organometallic-only,
        \texttt{compat(Pt,O)}$\!=\!0.90$ dative strong).

  \item \textbf{Valence-aware bond cap.}
        $\displaystyle \mathcal{L}_{\text{val}} =
        \sum_a \max\!\bigl(0,\,
        \text{cap}_a - \!\!\sum_{e\ni a}\!\!o_e + \varepsilon\bigr)^{\!2}$,
        per-atom cap $\text{cap}_a$ from
        \texttt{DEFAULT\_VALENCES} (C$=4$, N$=3$, O$=2$, F$=1$,
        Cl$=1$, Br$=1$, I$=1$; Pt$_II$=4 square planar,
        Cu$_II$=4, Au$_III$=4, Ru$_II$=6 octahedral,
        Ir$_III$=6). The squared-hinge form replaces the log
        barrier originally advertised in the design doc ---
        log barrier diverges too sharply for an untrained head;
        the squared hinge trades sharpness for stability
        (design-deviation documented in
        \texttt{decoder\_rework.py:419--465}).
\end{enumerate}

\paragraph{Combined prior and logit emission.}
$\displaystyle p_{\text{combined}} = p_{\text{dist}} \times p_{\text{type}}$
is mapped to a logit $\ell = \log(p/(1-p))\!\in\![-20,\,20]$ that
the inner \texttt{BondAwareDecoder}'s \texttt{BondOrderHead}
re-scores for the bond \emph{order} (SINGLE/DOUBLE/TRIPLE/AROMATIC).
The \texttt{ReworkedDecoder} decorator
(\texttt{molmetal/molmetal\_lam/lam\_chem/decoder\_rework.py:845})
composes the soft prior with the legacy decoder: the prior
filters candidate pairs to a "viable" set ($p_{\text{combined}}\!>\!0.5$),
then the inner decoder assigns bond orders. All three priors
are differentiable w.r.t.\ \texttt{coords} and the bond
\emph{logits}, so the joint PCGrad multi-task loss
\cite{yu2020pcgrad} can back-propagate through them in the
next retrain cycle.

\paragraph{Smoke validation (measured).}
On 192 synthetic CFM-style 10-atom clouds with realistic
1-5\,\AA{} coordinate spread (matches the round-12 trained-CFM
distribution), the soft prior lifts bond-bearing decode from
$0/192$ (Path A: legacy \texttt{BondAwareDecoder} with hard
2.4\,\AA{} cutoff) to $192/192$ (Path B: \texttt{ReworkedDecoder}
with soft 3-prior) at a 500-step + $h\!=\!64$ mini-budget.
96/192 of the Path-B mols fail \texttt{Chem.SanitizeMol} due
to overvalent atoms in the random cloud (e.g. C with 8 bonds)
--- this is a property of the synthetic coordinate generator,
not the decoder; the trained CFM will produce saner coords.
The key contract is that the rework \emph{proposes} bonds where
the legacy decoder proposes none; sanitization is downstream.

\paragraph{Relationship to Himo 2005 and the Pt-click matrix.}
The Himo 2005 CuAAC analysis identifies terminal alkyne +
azide dative bonding as the regioselectivity-determining step
for Pt-click scaffolds \cite{himo2005cuaac}. The
\texttt{ATOM\_TYPE\_COMPAT[(6, 7)]}\!=\!0.95 (C-N dative
strong) and \texttt{ATOM\_TYPE\_COMPAT[(17, 78)]}\!=\!0.95
(Pt-Cl cisplatin-style strong) cells directly encode these
two interactions. The Lit-Survey-v2 5$\times$5 Pt-click
compat matrix extends this to a 2-axis (donor atom-type
$\times$ Pt oxidation-state) compatibility grid used by the
scaffold-aware soft prior in
\texttt{r4\_lambda\_only\_run.py:236--267}.

% --- end-of-section pointer -------------------------------------------------
\paragraph{Cross-section pointers.}
Figure~\ref{fig:mlc-arch} (9-layer MLC architecture) is introduced
in \S\ref{sec:mlc-layers} and re-cited by \S\ref{sec:click-chem}
and \S\ref{sec:metal-geometry-prior}.
Figure~\ref{fig:click-reactions} (5 click reaction panels) is
introduced in \S\ref{sec:click-chem-smarts}.
Figure~\ref{fig:pipeline} (pocket-conditioned Lambda pipeline) is
cited by \S\ref{sec:metal-geom-enum} and
\S\ref{sec:mcts-closure-cross-ref}.
The closure theorem in Appendix~\ref{appendix:closure-theorem}
backs the strong-normalisation claim of
\S\ref{sec:mlc-strong-norm} and the sample-complexity claim of
\S\ref{sec:mcts-closure-cross-ref}; the proof is
\textsc{projected} (WF-Lambda-4, status: proof pending, machine
check deferred to supplementary material).
The secondary geometric generator (\S\ref{sec:method:secondary-cfm})
is \textsc{deferred} on this round and contributes no measurable
data cells; the deferred $\Lambda \times$ CFM coupling plan
(\texttt{TODO/pending/21\_lambda\_model\_coupling.md}) is
\emph{not} part of this section; it is mentioned in
\S\ref{sec:future} as future work.